% note.tex --- standalone math.OA / math-ph paper
% An Algebraic Weyl Curvature Hypothesis: Past Big Bang Non-Cyclic-Separating
% in FRW and the Full Bianchi Classification.
%
% Author: Kevin Remondiere
% Date: 2026-05-03
% Compile: pdflatex note ; bibtex note ; pdflatex note ; pdflatex note
%
% Builds on:
%   paper/frw_typeII_note/frw_note.tex      (Theorem 3.5, conformal pullback)
%   notes/path_realiste_2026_05_03/algebraic_arrow.{tex,py}   (FRW T1+T2+T3)
%   notes/path_realiste_2026_05_03/bianchi_extension_opus.{tex,py} (Prop. B.1)
%   notes/T2_bianchi_2026_05_03_jour2/T2_bianchi_extension.{tex,py} (T2-Bianchi I rigorous)
%   notes/T2_bianchi_2026_05_03_jour2/T2_bianchi_IX.{tex,py}        (T2-BIX pathwise)
%   notes/T2_bianchi_2026_05_03_jour2/T2_bianchi_V.{tex,py}         (T2-BV matter)
%   notes/T2_bianchi_2026_05_03_jour2/T2_bianchi_typeA.{tex,py}     (II, VI_0, VII_0, VIII)
%   notes/T2_bianchi_2026_05_03_jour2/T2_bianchi_typeB.{tex,py}     (III, IV, VI_h, VII_h)
%   notes/T2_bianchi_2026_05_03_jour2/hadamard_BIX_anisotropic.{tex,py}
%   notes/T2_bianchi_2026_05_03_jour2/hadamard_BV_anisotropic.{tex,py}

\documentclass[11pt]{article}

\usepackage[utf8]{inputenc}
\usepackage[T1]{fontenc}
\usepackage[margin=1in]{geometry}
\usepackage{amsmath, amssymb, amsthm}
\usepackage{mathtools}
\usepackage{hyperref}
\usepackage{xcolor}
\usepackage{cite}
\usepackage{booktabs}
\hypersetup{colorlinks=true, linkcolor=blue, citecolor=blue, urlcolor=blue}

% Theorem environments
\newtheorem{theorem}{Theorem}
\newtheorem{proposition}[theorem]{Proposition}
\newtheorem{lemma}[theorem]{Lemma}
\newtheorem{corollary}[theorem]{Corollary}
\theoremstyle{definition}
\newtheorem{definition}[theorem]{Definition}
\newtheorem{remark}[theorem]{Remark}
\newtheorem{question}[theorem]{Open Question}

% Macros (matching paper/frw_typeII_note/frw_note.tex style)
\newcommand{\R}{\mathbb{R}}
\newcommand{\C}{\mathbb{C}}
\newcommand{\Z}{\mathbb{Z}}
\newcommand{\HH}{\mathcal{H}}
\renewcommand{\AA}{\mathcal{A}}
\newcommand{\MM}{\mathcal{M}}
\newcommand{\NN}{\mathcal{N}}
\newcommand{\OO}{\mathcal{O}}
\newcommand{\Aut}{\operatorname{Aut}}
\newcommand{\Ad}{\operatorname{Ad}}
\newcommand{\BW}{\mathrm{BW}}
\newcommand{\HL}{\mathrm{HL}}
\newcommand{\FRW}{\mathrm{FRW}}
\newcommand{\Mink}{\mathrm{M}}
\newcommand{\BI}{\mathrm{B\text{-}I}}
\newcommand{\BII}{\mathrm{B\text{-}II}}
\newcommand{\BIII}{\mathrm{B\text{-}III}}
\newcommand{\BIV}{\mathrm{B\text{-}IV}}
\newcommand{\BV}{\mathrm{B\text{-}V}}
\newcommand{\BVIO}{\mathrm{B\text{-}VI_0}}
\newcommand{\BVIIO}{\mathrm{B\text{-}VII_0}}
\newcommand{\BVIh}{\mathrm{B\text{-}VI_h}}
\newcommand{\BVIIh}{\mathrm{B\text{-}VII_h}}
\newcommand{\BVIII}{\mathrm{B\text{-}VIII}}
\newcommand{\BIX}{\mathrm{B\text{-}IX}}
\newcommand{\BB}{\mathrm{BB}}
\newcommand{\HHH}{\mathbb{H}}
\newcommand{\Tr}{\operatorname{Tr}}
\DeclareMathOperator{\Sp}{Sp}
\DeclareMathOperator{\spec}{spec}

\title{An Algebraic Weyl Curvature Hypothesis:\\
Past Big Bang Non-Cyclic-Separating in FRW\\ and the Full Bianchi Classification}
\author{Kevin Remondi\`ere\\
\small Independent researcher\\
\small \href{mailto:kevin.remondiere@gmail.com}{kevin.remondiere@gmail.com}\\
\small ORCID: \href{https://orcid.org/0009-0008-2443-7166}{0009-0008-2443-7166}}
\date{2026-05-03}

\begin{document}
\maketitle

\begin{abstract}
We promote Penrose's Weyl Curvature Hypothesis (WCH) from a postulate on classical curvature at the past cosmological singularity to a theorem in algebraic quantum field theory (AQFT) on conformally flat FRW backgrounds and across the full nine-type Bianchi classification. The reformulation, which we call the \emph{Algebraic Weyl Curvature Hypothesis} (AWCH), is the statement that the inductive-limit local algebra of comoving causal diamonds extending to the Big Bang admits no Hadamard cyclic-separating vector, in contrast to the future-infinity inductive-limit local algebra which admits the Bunch--Davies vector and gives a type II$_\infty$ crossed product. For FRW we prove three independent statements: (T1) future infinity is type II$_\infty$ standard via Verch's inductive-limit theorem; (T2) the past Big Bang inductive-limit algebra has no Hadamard cyclic-separating vector, modulo the standard AQFT convention that physical states are Hadamard~\cite{HollandsWald2001,BrunettiFredenhagenVerch2003}; (T3) the algebraic past$\leftrightarrow$future asymmetry is intrinsic. The Bianchi extension is then carried out per type, with sharp labels: rigorous (Bianchi I), pathwise on the BKL attractor (Bianchi IX, via Heinzle--Uggla 2009~\cite{HeinzleUggla2009,HeinzleUgglaAttractor2009} and a Peter--Weyl SU(2) Hadamard SLE construction on Kasner-ordered trajectories), unconditional via Kontorovich--Lebedev decomposition (Bianchi V matter, with a leading-order Mehler--Sonine cancellation), rigorous-pending-Hadamard with explicit plausibility paths (Bianchi II, VI$_0$, VII$_0$, VIII), partial S3-conditional Hadamard-open (Bianchi III, IV, VI$_h$, VII$_h$), and \emph{partial S3-epoch-conditional} (Bianchi VI$_{-1/9}$, sharpened from a candidate clean negative via Lappicy--Uggla 2024~\cite{LappicyUggla2024} heteroclinic chains visiting Kasner-I vertices). The two universal driver obstructions are (S1)~a logarithmic-squared divergence of the volume-element zero-mode and (S3)~a long-wavelength tachyonic divergence along the contracting Kasner direction; both are intrinsically tied to non-vanishing Weyl curvature, propagated to the entire BFV-Verch local-quasi-equivalence folium~\cite{Verch1994,BrunettiFredenhagenVerch2003}. As a partial NO-GO byproduct, we identify a new \emph{open-FRW dichotomy}: in isotropic open ($k=-1$) FRW, the $H^3$ spectral gap removes the FRW zero-mode log divergence, leaving only the rescaling-failure obstruction; the algebraic past asymmetry in open FRW is therefore strictly weaker than in flat FRW, in contrast to Penrose's geometric WCH which treats both isotropic spatial topologies symmetrically. We compare with Penrose's geometric WCH: in FRW, $C_{abcd}\equiv 0$ and Penrose's content is trivial; in vacuum Kasner and in all genuinely anisotropic Bianchi types, the Weyl tensor diverges at the singularity and Penrose's WCH is geometrically violated, yet the algebraic WCH continues to hold (or to be conditionally established) wherever a Hadamard state has been constructed. Open extensions to non-conformal couplings, matter content beyond the strong-energy-condition asymptotic Kasner regime, and full quantum gravity beyond AQFT in fixed backgrounds are discussed.
\end{abstract}

\section{Introduction}\label{sec:intro}

Penrose's Weyl Curvature Hypothesis (WCH), formulated in the 1979 Einstein centenary volume~\cite{Penrose1979} and developed at length in \emph{Cycles of Time}~\cite{Penrose2010} (and in the conformal-cyclic-cosmology preprint arXiv:1011.3706), postulates that the Weyl curvature tensor $C_{abcd}$ vanishes (or is uniformly small in a suitable conformal sense) at the initial cosmological singularity. Together with non-vanishing Weyl elsewhere, this asymmetry selects a low-entropy past boundary condition and constitutes a possible foundation for the thermodynamic arrow of time.

Penrose's WCH is a \emph{geometric} hypothesis on the classical Weyl curvature at the singularity, not a derived consequence of dynamics. The present paper attempts to put it in a different form, derived from a hypothesis on the local algebras of quantum field theory (rather than from a hypothesis on classical curvature), within the algebraic-QFT framework of Brunetti--Fredenhagen--Verch~\cite{BrunettiFredenhagenVerch2003} and Hollands--Wald~\cite{HollandsWald2001}. Specifically, we replace ``Weyl tensor vanishes at past singularity'' with the operationally meaningful statement: \emph{the inductive-limit local algebra of comoving causal diamonds extending to the past singularity admits no cyclic-separating vector in any Hadamard GNS representation}.

This statement, which we call the \emph{Algebraic Weyl Curvature Hypothesis} (AWCH), is on the operator-algebraic level rather than the geometric level. It is automatic in the sense that, once formulated, it reduces to a theorem rather than a postulate, in two regimes that we treat below: spatially flat conformally coupled FRW (where Penrose's geometric content is trivial because $C_{abcd}\equiv 0$ throughout), and vacuum Bianchi I (where Penrose's geometric content is non-trivial and \emph{violated} in the sense $C_{txtx}\to\infty$ at the past). The algebraic statement therefore captures a different and complementary content from Penrose's: it remains a positive past$\leftrightarrow$future asymmetry of the local algebras even in regimes where Penrose's WCH is geometrically false.

\medskip

\noindent\textbf{Companion to.} The present paper builds on our companion technical note~\cite{frw_note} on the type II$_\infty$ classification of FRW comoving-diamond algebras via conformal pullback to Hislop--Longo modular flow, and on the post-CLPW operator-algebraic framework of Witten~\cite{Witten2022}, Chandrasekaran--Longo--Penington--Witten~\cite{CLPW2023}, and Faulkner--Speranza~\cite{FaulknerSperanza2024}. Sympy verifications of all numerical and symbolic intermediate steps are recorded in the source-bundle scripts \texttt{algebraic\_arrow.py}, \texttt{bianchi\_extension\_opus.py}, \texttt{T2\_bianchi\_extension.py}.

\medskip

\noindent\textbf{Paper structure.} Section~\ref{sec:prelim} gathers preliminaries on Tomita--Takesaki, BFV local quasi-equivalence, Hadamard states, and the modular crossed product. Section~\ref{sec:frw} restates the FRW theorem T1+T2+T3 from~\cite{algebraic_arrow}. Section~\ref{sec:bianchi-T1} restates the perturbative Bianchi I T1 lift (Proposition B.1 of~\cite{bianchi_extension_opus}). Section~\ref{sec:bianchi-T2} contains the rigorous Theorem T2-Bianchi I, with explicit proof and sympy verifications. Section~\ref{sec:bianchi-IX} extends T2 to vacuum Bianchi IX pathwise on the BKL Kasner attractor via the Heinzle--Uggla 2009 attractor theorem~\cite{HeinzleUgglaAttractor2009}, complemented by a Peter--Weyl SU(2) Hadamard SLE construction valid on Kasner-ordered trajectories. Section~\ref{sec:bianchi-V} gives an unconditional T2 for matter Bianchi V via the Kontorovich--Lebedev decomposition on $H^3$ (Lebedev 1965~\cite{Lebedev1965}, Helgason 2008~\cite{Helgason2008}) with a Mehler--Sonine UV cancellation, and a leading-order Hadamard SLE construction. Section~\ref{sec:typeA-other} treats the remaining Type A signatures (II, VI$_0$, VII$_0$, VIII) per type with S1+S3 status table and Hadamard plausibility. Section~\ref{sec:typeB-other} treats Type B (III, IV, VI$_h$, VII$_h$) and the partial S3-epoch-conditional verdict for VI$_{-1/9}$ via the Lappicy--Uggla 2024 heteroclinic-chain refinement~\cite{LappicyUggla2024} of the Hewitt--Horwood--Wainwright 2003 exceptional locus~\cite{HewittHorwoodWainwright2003}. Section~\ref{sec:openFRW} establishes the partial NO-GO open-FRW dichotomy: the $H^3$ spectral gap kills the FRW zero-mode log divergence, leaving T2 strictly weaker in open FRW than in flat FRW. Section~\ref{sec:penrose} compares with Penrose's geometric WCH. Section~\ref{sec:conclusion} concludes with a unified table of all nine Bianchi types and a brief CCM-Bianchi IX bridge teaser~\cite{CCM_BIX_dictionary_max}.

\medskip

\noindent\textbf{Honest framing of conditional vs.~unconditional.} The status labels used throughout are: \emph{rigorous} (Bianchi I, Theorem~\ref{thm:T2-BI}, conditional only on the Banerjee--Niedermaier~\cite{BanerjeeNiedermaier2023} Hadamard SLE on $T^3$/$\R^3$); \emph{pathwise on the BKL attractor} (Bianchi IX, Theorem~\ref{thm:T2BIX}, conditional on Hadamard quasifree existence which we partially supply on Kasner-ordered trajectories); \emph{unconditional modulo a 4--6-week mechanical write-up of one Olver-style WKB lemma} (matter Bianchi V, Theorem~\ref{thm:T2BV}); \emph{rigorous pending Hadamard} (Bianchi II, VI$_0$, VII$_0$, VIII, Theorems~\ref{thm:T2-BII}--\ref{thm:T2-BVIII}); \emph{partial / S3-conditional, Hadamard-open} (Type B III, IV, VI$_h$, VII$_h$, Theorem~\ref{thm:T2-typeB}); \emph{partial / S3-epoch-conditional} (Bianchi VI$_{-1/9}$, Proposition~\ref{prop:T2-VI19blocked}, sharpened from the prior ``BLOCKED'' candidate via Lappicy--Uggla 2024~\cite{LappicyUggla2024} heteroclinic chains). All Hadamard convention statements are within the standard AQFT convention~\cite{HollandsWald2001,BrunettiFredenhagenVerch2003,Verch1994} that physical states are Hadamard. We make no claim of unconditional rigour outside these explicitly stated convention scopes.

\section{Preliminaries}\label{sec:prelim}

We work throughout on globally hyperbolic spacetimes $(M,g)$ in $d=4$ with the conformally coupled massless scalar field $\varphi$ ($\xi=1/6$, $m=0$), whose classical action is~\cite[\S4.6]{Wald1994}
\begin{equation}\label{eq:S-cc}
S[\varphi] = -\tfrac{1}{2}\int \mathrm{d}^4 x\, \sqrt{-g}\,\Big[g^{\mu\nu}\partial_\mu\varphi\,\partial_\nu\varphi - \tfrac{1}{6}\,R\,\varphi^2\Big],
\end{equation}
with EOM $(\Box_g - R/6)\varphi=0$ and conformal-covariance identity $\Box_{\tilde g}(a\varphi) = a^3(\Box_g - R/6)\varphi$ under $g_{\mu\nu}=a^2\tilde g_{\mu\nu}$~\cite[\S4.6]{Wald1994}.

\subsection{Local algebras and Tomita--Takesaki}\label{sec:tomita}

To each open relatively compact region $\OO\subset M$ with causally complete closure we associate the smeared local field algebra
\begin{equation}\label{eq:algebra-def}
\AA(\OO) := \pi_\omega\bigl(\{\varphi(f) : f\in C_c^\infty(\OO)\}\bigr)'',
\end{equation}
the double commutant in the GNS representation $\pi_\omega$ of a fixed Hadamard reference state $\omega$ on the field algebra of $(M,g)$. Standard Haag--Kastler axioms (isotony, locality, covariance) hold; existence of a Hadamard state on globally hyperbolic spacetimes is guaranteed by the deformation argument of Fulling--Narcowich--Wald and refined by Junker, with constructive families on FRW (the Bunch--Davies / conformal vacuum~\cite[\S5.4]{Wald1994}) and Bianchi I (Banerjee--Niedermaier~\cite{BanerjeeNiedermaier2023}; Brum--Them~\cite{BrumThem2013} for inhomogeneous spacetimes).

If a vector $\Omega\in\HH_\omega$ is cyclic and separating for $\AA(\OO)$, Tomita--Takesaki theory associates to the pair $(\AA(\OO),\Omega)$ the modular operator $\Delta$, modular conjugation $J$, and the modular automorphism group $\sigma_t = \mathrm{Ad}(\Delta^{it})$~\cite{BratteliRobinsonII}. The Connes--Takesaki classification of type III factors~\cite{Connes1973} attaches the spectral invariant $\Gamma(\sigma)$ (Connes spectrum); when $\Gamma(\sigma)=\R$ and $\AA(\OO)$ is type III$_1$, the crossed product $\AA(\OO)\rtimes_\sigma\R$ is a type II$_\infty$ factor~\cite{TakesakiII}.

\subsection{Brunetti--Fredenhagen--Verch local quasi-equivalence}\label{sec:bfv}

The Brunetti--Fredenhagen--Verch generally covariant locality principle~\cite{BrunettiFredenhagenVerch2003} casts AQFT on curved spacetime as a covariant functor from the category of globally hyperbolic spacetimes (with morphisms isometric embeddings preserving causal structure) to the category of unital $C^*$-algebras (with morphisms unital injective $*$-homomorphisms). A central technical consequence, exploited throughout the present paper, is the \emph{local quasi-equivalence} of any two Hadamard quasifree states on the same algebra: their GNS representations are quasi-equivalent on every relatively compact diamond bounded away from singularities, with the equivalence implemented (locally) by a unitary isomorphism preserving the Wightman positivity structure~\cite{Verch1994}.

\subsection{Hadamard states and the standard AQFT convention}\label{sec:hadamard}

A quasifree state $\omega$ on the field algebra of $(M,g)$ is \emph{Hadamard} if its two-point function $W_\omega(x,y) = \langle\varphi(x)\varphi(y)\rangle_\omega$ has the locally specified short-distance behaviour
\begin{equation}\label{eq:hadamard}
W_\omega(x,y) = \frac{U(x,y)}{8\pi^2 \sigma_+(x,y)} + V(x,y)\log\sigma_+(x,y) + W_{\mathrm{smooth}}(x,y),
\end{equation}
with $U$, $V$ Hadamard parametrices determined by the local geometry and EOM, $\sigma_+$ the Synge world function with the standard $i\epsilon$ prescription, and $W_{\mathrm{smooth}}$ a state-dependent smooth bi-distribution~\cite{HollandsWald2001}. Two Hadamard states $\omega,\omega'$ on the same $(M,g)$ have the same $U$ and $V$, and their difference $W_\omega - W_{\omega'}$ is a smooth bi-distribution. \emph{Standard convention} of Hollands--Wald and BFV: physical states are Hadamard. The present paper adopts this convention, and theorem statements that reference ``no Hadamard cyclic-separating vector'' are within this scope.

\subsection{Modular crossed product and CLPW type II classification}\label{sec:crossed}

Given a type III$_1$ von Neumann algebra $\MM$ and the modular automorphism group $\sigma:\R\to\Aut(\MM)$ of a faithful normal cyclic-separating state, the crossed product $\MM\rtimes_\sigma\R$ on $L^2(\R,\HH)$ is generated by $\pi(m)\xi(s)=\sigma_{-s}(m)\xi(s)$ and translation $u_t\xi(s)=\xi(s-t)$~\cite{TakesakiII,Connes1973}. By Connes--Takesaki, $\MM\rtimes_\sigma\R$ is a type II$_\infty$ factor when $\Gamma(\sigma)=\R$ (modular Hamiltonian unbounded both ways)~\cite{Witten2022,CLPW2023,FaulknerSperanza2024}. The CLPW dichotomy distinguishes type II$_\infty$ (unbounded modular Hamiltonian, Minkowski wedge / Hislop--Longo diamond / Faulkner--Speranza horizon settings) from type II$_1$ (bounded-below modular Hamiltonian, de Sitter static patch).

\section{FRW Algebraic WCH}\label{sec:frw}

We restate the three FRW theorems established in our companion technical note~\cite{algebraic_arrow}, which builds on Theorem 3.5 of~\cite{frw_note}. We work throughout on radiation-dominated FRW with $a(\eta)=a_0\eta$ on $\eta\in(0,\infty)$ in conformal time; the matter-dominated case $a(\eta)\propto\eta^2$ is treated identically~\cite[Theorem 3.6]{frw_note}.

Let $D_n$ denote the increasing family of comoving causal diamonds with fixed $\eta_i>0$ and $\eta_f^{(n)}=n\to+\infty$, and let $D_\BB^{(n)}$ denote the decreasing family with fixed $\eta_f<\infty$ and $\eta_i^{(n)}\to 0^+$. Define the inductive-limit algebras
\begin{equation}\label{eq:Dlims-frw}
\AA(D_\infty)_\FRW := \overline{\bigcup_n \AA(D_n)_\FRW}^{\mathrm{w}^*},
\qquad
\AA(D_\BB)_\FRW := \overline{\bigcup_n \AA(D_\BB^{(n)})_\FRW}^{\mathrm{w}^*}.
\end{equation}

\begin{theorem}[T1 -- future-infinity II$_\infty$ standard]\label{thm:T1}
For the conformally coupled massless scalar on radiation-dominated FRW in the conformal vacuum $\Omega_\FRW$: $\AA(D_\infty)_\FRW$ is a well-defined W$^*$-algebra; $\Omega_\FRW$ extends to a cyclic-separating vector for $\AA(D_\infty)_\FRW$ in the inductive-limit GNS representation; and the crossed product $\AA(D_\infty)_\FRW\rtimes_{\sigma^\FRW}\R$ is a type II$_\infty$ factor.
\end{theorem}

\begin{proof}[Proof sketch]
For each finite $D_n$, the conformal-rescaling unitary $U_n$ of~\cite[Prop.~3.3]{frw_note} gives $\Ad U_n:\AA(D_n)_\FRW\xrightarrow{\sim}\AA(M_n)_\Mink$ with $U_n\Omega_\FRW=\Omega_\Mink$. The $\{U_n\}$ are coherent on overlaps because $U_n|_{\AA(D_m)_\FRW}=U_m$ for $m\le n$, and inductively yield a unitary $U_\infty:\HH_\FRW\to\HH_\Mink$ implementing $\AA(D_\infty)_\FRW\simeq\AA(M_\infty)_\Mink$ with $U_\infty\Omega_\FRW=\Omega_\Mink$. The Minkowski limit $\AA(M_\infty)_\Mink$ is a standard local algebra in Minkowski AQFT; the Minkowski vacuum $\Omega_\Mink$ remains cyclic and separating for the limit, by Verch's inductive-limit theorem on Hadamard vacuum representations~\cite{Verch1997}. Connes--Takesaki then gives the type II$_\infty$ classification of the crossed product, with full Connes spectrum $\Gamma(\sigma)=\R$ inherited from each finite stage by Hislop--Longo~\cite{HislopLongo1982}.
\end{proof}

\begin{theorem}[T2 -- past Big Bang non-cyclic-separating]\label{thm:T2-FRW}
For the same theory, $\Omega_\FRW$ does not extend to a cyclic-separating vector for $\AA(D_\BB)_\FRW$, and no Hadamard state on $(M,g_\FRW)$ in the BFV-Verch local-quasi-equivalence folium of $\Omega_\FRW$ extends to one. Equivalently, the smeared two-point function is unbounded on the natural test-function space $C_c^\infty(D_\BB)$ extending to $\eta=0$, and no positive bilinear form on this space yields a bounded GNS representation.
\end{theorem}

\begin{proof}[Proof sketch]
The proof has three sympy-verifiable obstructions, each of which alone suffices to obstruct cyclic-separation; we list them by their CHECK label in \texttt{algebraic\_arrow.py}.

\textbf{(i) Rescaling failure (CHECK 2).} The map $f\mapsto a^{-3}f$ is required for the conformal-pullback unitary to invert. For $h(\eta)=\eta^k/k!$ generic Minkowski test function, $h/a^3 = \eta^{k-3}/(k!\,a_0^3)$ diverges at $\eta=0$ for $k=0,1,2$. Generic $h\in C_c^\infty(M_\BB)$ has nonzero Taylor coefficients $c_0,c_1,c_2$ at $\eta=0$, hence is not in the image of any $f\in C_c^\infty(D_\BB)$ under $f\mapsto a^3 f$. The Prop.~3.3 isomorphism of~\cite{frw_note} fails on $D_\BB$.

\textbf{(ii) IR divergence of smeared two-point (CHECK 3).} Off-lightcone with spatial separation $r>0$ fixed,
\begin{equation*}
\langle\varphi(f)^2\rangle_{\Omega_\FRW}
\;\ge\; C(r)\,\biggl|\!\int\!\!\int_\delta^{\delta+\epsilon} \frac{f(\eta_x)f(\eta_y)}{a_0^2\,\eta_x\,\eta_y}\,\mathrm{d}\eta_x\mathrm{d}\eta_y\biggr|
\;=\; C(r)\,\frac{\log^2(\delta/(\delta+\epsilon))}{a_0^2}\;\xrightarrow[\delta\to 0^+]{}\;+\infty,
\end{equation*}
i.e.\ the limiting bilinear form is unbounded. The single-integral version $\int_0^\epsilon\mathrm{d}\eta/\eta = +\infty$ is sharp.

\textbf{(iii) Hadamard exclusion.} Any state $\omega$ in the Hadamard folium of $\Omega_\FRW$ has $W_\omega - W_{\Omega_\FRW}$ smooth on $D_\BB\times D_\BB$ by~\eqref{eq:hadamard} and~\cite{HollandsWald2001,BrunettiFredenhagenVerch2003}. The smooth correction cannot cure the singular prefactor $1/(a(\eta_x)a(\eta_y))$. Hence (ii) extends to all Hadamard-equivalent states.

A cyclic-separating vector for $\AA(D_\BB)_\FRW$ in any Hadamard GNS representation would yield, via Tomita--Takesaki, a finite modular operator and hence a bounded bilinear form, contradicting (iii). Hence no such vector exists.
\end{proof}

\begin{remark}[Hadamard convention scope]\label{rem:hadamard-scope}
Step (iii) covers the Hadamard folium of $\Omega_\FRW$: states quasi-equivalent to the conformal vacuum. The standard AQFT-on-curved-spacetime convention is that physical states are Hadamard (Section~\ref{sec:hadamard}); under this convention, T2 is unconditional. A residual theoretical possibility (excluded by the convention) is that some non-Hadamard state with non-Hadamard high-frequency tails extends to a cyclic-separating vector. We adopt the standard convention and treat this possibility as outside scope.
\end{remark}

\begin{theorem}[T3 -- algebraic arrow of time]\label{thm:T3}
The W$^*$-algebraic pairs $(\AA(D_\infty)_\FRW,\Omega_\FRW)$ and $(\AA(D_\BB)_\FRW,\Omega_\FRW)$ are not isomorphic, and there is no W$^*$-isomorphism between them intertwining the time-reversal $\eta\mapsto-\eta$ within the FRW chart. The asymmetry between past Big Bang ($\eta\to 0^+$) and future infinity ($\eta\to+\infty$) is therefore an intrinsic invariant of the local-algebra net of $(M,g_\FRW)$, not an additional postulate.
\end{theorem}

\begin{proof}
By Theorem~\ref{thm:T1}, $\AA(D_\infty)_\FRW$ admits a cyclic-separating vector. By Theorem~\ref{thm:T2-FRW}, $\AA(D_\BB)_\FRW$ does not (within the Hadamard folium). A W$^*$-isomorphism between the pairs would have to map cyclic-separating to cyclic-separating, since this is a W$^*$-algebraic invariant of the (algebra, state) pair: hence no such isomorphism exists. The time-reversal $\eta\mapsto T_*-\eta$ for any finite $T_*$ yields no cure: there is no FRW-internal singularity at finite $\eta>0$ along a comoving worldline (the FRW chart is geodesically complete to the future for radiation FRW, with $a(\eta)\to+\infty$). For matter-dominated FRW the same argument applies, with strictly stronger past obstruction $\int_0^\epsilon\mathrm{d}\eta/\eta^2=+\infty$ (CHECK 5).
\end{proof}

\section{Bianchi I T1 perturbative lift}\label{sec:bianchi-T1}

We restate the perturbative lift of T1 from FRW to small-anisotropy Bianchi I, originally Proposition B.1 of our companion note~\cite{bianchi_extension_opus}. The proof is by Connes 1973 cocycle invariance of the Connes spectrum~\cite{Connes1973} together with BFV local quasi-equivalence~\cite{BrunettiFredenhagenVerch2003,Verch1994}.

\begin{proposition}[Perturbative type II$_\infty$ for near-FRW Bianchi I]\label{prop:B1}
Let $\bigl(M, g_\BI(\varepsilon)\bigr)$ be a one-parameter family of Bianchi I spacetimes with $g_\BI(0) = g_\FRW$ (radiation- or matter-dominated, spatially flat $k=0$) and small anisotropy parameter $\varepsilon\in[0,\varepsilon_0)$:
\begin{equation}\label{eq:perturb-metric}
\mathrm{d}s^2_{\BI(\varepsilon)} \;=\; -\mathrm dt^2
   \,+\, \sum_{i=1}^3 a(t)^2 \bigl(1 + \varepsilon\,\delta_i(t)\bigr)^2\,\mathrm dx^i\,\mathrm dx^i,
\qquad \sum_i \delta_i \equiv 0
\end{equation}
(traceless anisotropy, volume preserving to first order). Let $\varphi$ be the conformally coupled massless scalar in $d=4$, and $\Omega_{\BI(\varepsilon)}$ the SLE / states-of-low-energy ground state of \cite{BanerjeeNiedermaier2023, BrumThem2013}. Let $D\subset M$ be a doubly-bounded comoving causal diamond with $0<t_i<t_f<\infty$. Then there exists $\varepsilon_0 = \varepsilon_0(D)>0$ such that for all $\varepsilon\in[0,\varepsilon_0)$:
\begin{enumerate}
\item $\AA(D)_{\BI(\varepsilon)}$ is a type III$_1$ factor in the GNS rep of $\Omega_{\BI(\varepsilon)}$.
\item $\Gamma\bigl(\sigma^{\BI(\varepsilon)}\bigr) = \R$.
\item $\AA(D)_{\BI(\varepsilon)} \rtimes_\sigma \R$ is a type II$_\infty$ factor.
\end{enumerate}
\end{proposition}

\begin{proof}
For $\varepsilon=0$ this is Theorem~3.5 of~\cite{frw_note} via the conformal-pullback unitary $U$.

\textbf{Step 1 (BFV local quasi-equivalence).} By Brunetti--Fredenhagen--Verch~\cite[Thm.~4.1]{BrunettiFredenhagenVerch2003} together with Hollands--Wald~\cite[\S5]{HollandsWald2001} and Verch~\cite{Verch1994}, the GNS representations $\pi_0:=\pi_{\Omega_\FRW}$ and $\pi_1:=\pi_{\Omega_{\BI(\varepsilon)}}$ on the local algebra $\AA(D)$ are quasi-equivalent: there exists a $W^*$-isomorphism $\pi_1(\AA(D))''\xrightarrow{\sim}\pi_0(\AA(D))''$ on the same closure $\bar\AA(D)$. The threshold $\varepsilon_0(D)$ is the BFV quasi-equivalence radius, estimated by $\varepsilon_0\sim (\text{curvature radius of }D)/(\text{diameter of }D)$, well within CMB-anisotropy bounds $\delta T/T\sim 10^{-5}$ on largest scales for cosmologically realistic $D$.

\textbf{Step 2 (Connes cocycle).} By Connes~\cite[Th\'eor\`eme 1.2.4]{Connes1973} (Radon--Nikodym for modular automorphisms): for faithful normal states $\omega_0,\omega_1$ on the same vN algebra $M$, the modular flows differ by a $\sigma$-cocycle $u_t := (D\omega_1\!:\!D\omega_0)_t$ inside $M$, hence
\begin{equation}\label{eq:connes-cocycle}
\sigma_t^{\omega_1} \;=\; \Ad u_t \circ \sigma_t^{\omega_0}.
\end{equation}
Apply with $M = \pi_0(\AA(D))''$, $\omega_0 = $ Bunch--Davies state on FRW, $\omega_1 = $ adiabatic Hadamard state on $\BI(\varepsilon)$ pushed through the BFV isomorphism of Step 1.

\textbf{Step 3 (Connes spectrum invariance).} By Connes~\cite[Cor.~1.2.4]{Connes1973}, the Connes spectrum is invariant under inner perturbation by a $\sigma$-cocycle inside $M$:
\begin{equation*}
\Gamma(\sigma^{\omega_1}) \;=\; \Gamma(\sigma^{\omega_0}) \;=\; \Gamma(\sigma^\HL) \;=\; \R,
\end{equation*}
the last equality by Hislop--Longo~\cite{HislopLongo1982}.

\textbf{Step 4 (type classification).} By Connes--Takesaki, type III$_1$ crossed by its modular automorphism group with full Connes spectrum $\Gamma=\R$ is type II$_\infty$.

The type III$_1$ factoriality of $\AA(D)_{\BI(\varepsilon)}$ (item~1) is preserved by the $W^*$-isomorphism of Step~1 from the FRW type III$_1$ factor $\AA(D)_\FRW$ (which is type III$_1$ by~\cite[Thm.~3.5]{frw_note} via Hislop--Longo).
\end{proof}

\noindent\textbf{Sympy verification (\texttt{bianchi\_extension\_opus.py} CHECK C4).} For radiation FRW $a(t)=\sqrt t$ and the perturbed metric~\eqref{eq:perturb-metric}, the Riemann component $R^t{}_{xtx}$ is computed as a Taylor series in $\varepsilon$:
\begin{equation*}
R^t{}_{xtx}\bigl(\varepsilon\bigr)
   \;=\; -\frac{1}{4t} \;-\; \frac{\delta_1\,\varepsilon}{2t}\;+\;O(\varepsilon^2),
\end{equation*}
with the $\varepsilon^0$ term equal to the FRW Riemann (Weyl-vanishing part) and the $\varepsilon^1$ term carrying the leading anisotropic Weyl correction. The expansion is well-defined and confirms $g_\BI(\varepsilon)$ is a smooth one-parameter family with $g_\BI(0)=g_\FRW$.

\section{Bianchi I T2 rigorous (new result)}\label{sec:bianchi-T2}

This section contains the central new result of the present paper: a rigorous theorem T2-Bianchi I, conditional only on the Banerjee--Niedermaier~\cite{BanerjeeNiedermaier2023} construction of Hadamard SLE states on Bianchi I.

\subsection{Statement and standing assumptions}\label{sec:T2-statement}

We use synchronous time $t>0$ throughout, with Big Bang at $t=0^+$. Vacuum Bianchi I (Kasner) has metric
\begin{equation}\label{eq:kasner}
\mathrm{d}s^2_{\BI} = -\mathrm{d}t^2 + \sum_{i=1}^3 t^{2 p_i}\,\mathrm{d}x_i^2,\qquad
\sum_i p_i = \sum_i p_i^2 = 1
\end{equation}
(vacuum Kasner constraint, $\sum_i p_i = 1$ coming from $T_{\mu\nu}=0$ and the $\sum_i p_i^2 = 1$ coming from the Hamiltonian constraint). The standard non-degenerate exponent triple is $(p_1,p_2,p_3)=(-1/3,2/3,2/3)$ (one contracting direction, two expanding), which is the only vacuum-Kasner solution up to permutations and the trivial $(1,0,0)$ branch (locally flat). Comoving causal diamonds are
\begin{equation}\label{eq:diamond-bi}
D_{(t_i,t_f)} := \{\,(t,\vec x)\,:\,t_i \le t \le t_f,\;|\vec x|_{g_3(t_i)} \le R\,\},
\end{equation}
and we form the inductive limit over $t_i\downarrow 0$ (at fixed $t_f$, $R$):
\begin{equation}\label{eq:Dlim-bi}
\AA(D_\BB)_\BI := \overline{\bigcup_{t_i\downarrow 0} \AA(D_{(t_i,t_f)})_\BI}^{C^*}.
\end{equation}

\begin{theorem}[T2-Bianchi I, vacuum case]\label{thm:T2-BI}
Let $(M,g)$ be vacuum Bianchi I with non-degenerate Kasner exponents $(p_1,p_2,p_3)=(-1/3,2/3,2/3)$ (or any permutation). Let $\varphi$ be the conformally coupled massless scalar field on $(M,g)$, and let $\omega$ be the SLE Hadamard quasifree state of Banerjee--Niedermaier~\cite{BanerjeeNiedermaier2023}.\footnote{Banerjee--Niedermaier construct their SLE on a $T^3$ spatial slice with periodic identification. The construction extends to $\R^3$ by a mode-decomposition argument (compactly supported test functions on $\R^3$ have well-defined Fourier transforms decaying at infinity). Existence of a Hadamard state on non-compact Bianchi I is also covered by~\cite{BrumThem2013} via the inhomogeneous-spacetime SLE construction.} Then the inductive limit~\eqref{eq:Dlim-bi} does not admit $\omega$, nor any state in the BFV-Verch local-quasi-equivalence folium of $\omega$, as a cyclic-separating vector in any GNS representation.
\end{theorem}

\subsection{Proof}\label{sec:T2-proof}

The proof is by establishing two independent infrared divergences of the smeared two-point function on test functions extending to $t=0$, then invoking BFV-Verch local quasi-equivalence to cover the entire Hadamard folium.

\paragraph{(i) Volume-element divergence (S1).}

For vacuum Bianchi I the spatial-volume element is
\begin{equation}\label{eq:vol-element}
\sqrt{g_3(t)} = a_1(t)\,a_2(t)\,a_3(t) = t^{p_1+p_2+p_3} = t,
\end{equation}
since the vacuum Kasner constraint $\sum_i p_i = 1$ enforces a linear vanishing of the volume. The state $\omega$, being Hadamard, has a Wightman two-point function whose state-dependent smooth part contains the $1/\sqrt{g_3(t_x)g_3(t_y)} = 1/\sqrt{t_x t_y}$ density factor (the standard mode-functional normalisation of any Hadamard state on a Bianchi I background; see~\cite[\S3]{BanerjeeNiedermaier2023} for the SLE form).

The Klein--Gordon zero-mode (constant in $\vec x$) component $T_0(t)$ of the field on vacuum Bianchi I satisfies (with $\sum p_i=1$)
\begin{equation}\label{eq:zero-mode-eq}
\ddot T_0 + \frac{1}{t}\,\dot T_0 = 0
\;\Longrightarrow\;
T_0(t) = c_1\log t + c_2,
\end{equation}
hence the smeared zero-mode contribution to $\langle\varphi(f)^2\rangle_\omega$ contains
\begin{equation}\label{eq:zero-mode-log-sq}
\int_\delta^{2\delta}\!\!\!\int_\delta^{2\delta}\!
   \frac{(\log t_x)\,(\log t_y)}{t_x\,t_y}\,\mathrm{d}t_x\,\mathrm{d}t_y
  \;=\; \biggl(\int_\delta^{2\delta}\!\frac{\log t}{t}\,\mathrm{d}t\biggr)^{\!\!2}
  \;=\; \biggl(\frac{(\log 2\delta)^2-(\log\delta)^2}{2}\biggr)^{\!\!2}
  \;\xrightarrow[\delta\to 0^+]{}\; +\infty.
\end{equation}
(Sympy-verified, see \texttt{T2\_bianchi\_extension.py} block S1.) Equation~\eqref{eq:zero-mode-log-sq} is the Bianchi I analog of the FRW prefactor divergence~\cite[CHECK~3]{algebraic_arrow}: the FRW $1/(a_0^2\eta_x\eta_y)$ factor is replaced by the $1/(t_x t_y)$ vacuum-Kasner zero-mode factor, with the same $\log^2$-divergence structure.

\paragraph{(ii) Long-wavelength tachyonic mode along the contracting Kasner direction (S3, NEW for Bianchi I).}

In addition to the FRW-analog divergence in (i), the genuine anisotropy of vacuum Bianchi I produces a SECOND, INDEPENDENT IR divergence absent in FRW. For modes with comoving wave-vector $\vec k=(k_1,0,0)$ aligned with the contracting direction $p_1=-1/3$, the local frequency is
\begin{equation}\label{eq:tachyonic-omega}
\omega_{\vec k}(t) = \sqrt{k_i k_j g_3^{ij}(t)}
   = |k_1|\,t^{-p_1} = |k_1|\,t^{1/3},
\end{equation}
so $\omega_{\vec k}(t)\to 0$ as $t\to 0^+$ for every fixed $k_1\neq 0$: the mode is \emph{effectively massless} at the singularity, behaving tachyonically (zero-frequency limit). The Banerjee--Niedermaier SLE mode functions~\cite[Theorem 4.3]{BanerjeeNiedermaier2023} have leading-order WKB form
\begin{equation}\label{eq:wkb-mode}
T_{\vec k}(t) \sim \frac{1}{\sqrt{2\,\omega_{\vec k}(t)\,\sqrt{g_3(t)}}}
  \;=\; \frac{1}{\sqrt{2|k_1|\,t^{1/3}\cdot t}}
   \;=\; \frac{1}{\sqrt{2|k_1|\cdot t^{4/3}}}.
\end{equation}
A $C_c^\infty$ test function $f(t,\vec x) = \chi_{[\delta,2\delta]}(t)\,h(\vec x)$ with $\hat h(\vec 0)\neq 0$ then produces a $k_1$-IR divergence in the smeared Wightman two-point:
\begin{align}
\langle\varphi(f)^2\rangle_\omega
  &\supset
  \int_{-K}^K \mathrm{d}k_1\,|\hat h(k_1,0,0)|^2\,
   \biggl|\int_\delta^{2\delta}T_{\vec k}(t)\,\mathrm{d}t\biggr|^2
  \notag\\
  &\sim |\hat h(0)|^2 \int_{-K}^K \frac{\mathrm{d}k_1}{|k_1|}\cdot O(\delta^{2/3})
  \notag\\
  &= +\infty\quad(\text{any }\delta>0,\text{ at }K\text{ fixed}).
  \label{eq:k-IR}
\end{align}
The IR divergence in $k_1$ is logarithmic and \emph{independent} of the time-cutoff $\delta$. (Sympy-verified, \texttt{T2\_bianchi\_extension.py} block S3.)

This is a phenomenon \emph{absent in FRW}: there, the conformal-vacuum mode functions $T_k^\FRW(\eta) = e^{-ik\eta}/(a_0\eta\sqrt{2k})$ are plane-wave bounded for every $\eta>0$, with no $k_1\to 0$ enhancement. The contracting Kasner direction is the algebraic source of the new divergence, and is intrinsically tied to the Weyl-non-vanishing character of vacuum Bianchi I.

\paragraph{(iii) BFV folium uniformity.}

Either of (i) or (ii) alone gives an IR divergence of $\langle\varphi(f)^2\rangle_\omega$ as $\delta\to 0^+$ for the SLE state $\omega$. By Verch's local quasi-equivalence theorem~\cite{Verch1994}, any other Hadamard quasifree state $\omega'$ on Bianchi I is locally quasi-equivalent to $\omega$ on every $\AA(D)$ with $D$ a relatively compact diamond bounded away from $t=0$. Equivalence on the inductive limit follows from the Brunetti--Fredenhagen--Verch generally covariant local quasi-equivalence~\cite{BrunettiFredenhagenVerch2003}: the difference of two-point functions is a smooth integral kernel, hence cannot cancel the $1/\sqrt{t_x t_y}$ singular prefactor or the $1/|k_1|$ tachyonic IR divergence. Therefore the divergence holds for every state in the BFV folium of $\omega$.

\paragraph{(iv) Conclusion.}

A cyclic-separating vector for $\AA(D_\BB)_\BI$ in the GNS representation of any Hadamard state in the BFV folium of $\omega$ would carry a normal state on the inductive-limit algebra. Such a state would have a finite smeared two-point on every test function in $C_c^\infty(D_\BB)$, including those with support extending to $t=0$. This contradicts (i), (ii), (iii). \qed

\subsection{Sympy verifications (summary)}\label{sec:T2-sympy}

We summarise the explicit sympy verifications underlying Theorem~\ref{thm:T2-BI}; the full script is \texttt{T2\_bianchi\_extension.py}.

\begin{itemize}
\item \textbf{S1 (volume zero-mode log$^2$).} Sympy computes $\int_\delta^{2\delta}\log t/t\,\mathrm{d}t = \tfrac12[(\log 2\delta)^2-(\log\delta)^2]$, then squares, then takes $\delta\to 0^+$, returning $+\infty$.
\item \textbf{C1 (Kasner Weyl divergence).} Sympy computes the Weyl tensor of the metric~\eqref{eq:kasner} with $(p_1,p_2,p_3)=(-1/3,2/3,2/3)$ and returns $C_{txtx}=-4/(9 t^{8/3})$, divergent at $t=0$.
\item \textbf{S3 (contracting-direction tachyon).} Sympy verifies that $\omega_{(k_1,0,0)}(t)=|k_1|\,t^{1/3}\to 0$ as $t\to 0^+$, that the WKB mode amplitude scales as $1/\sqrt{t^{4/3}}$, and that $\int dk_1/|k_1|=+\infty$ at any fixed UV cutoff $K$.
\item \textbf{C3 (Bianchi V Weyl).} Sympy returns $C_{txtx}=31/(864\,t)$ for the standard Bianchi V metric, divergent at $t=0$.
\item \textbf{Bianchi IX volume.} Sympy verifies the time-averaged Hubble rate on the Kasner-attractor invariant measure is bounded below by a positive constant; the volume $V(t)=a_1 a_2 a_3$ is monotonically vanishing $V(t)\to 0^+$ even through the chaotic Kasner-epoch sequence (Misner $\alpha\to-\infty$).
\end{itemize}

\subsection{Forward references to Bianchi V and Bianchi IX}\label{sec:T2-coverage}

The Bianchi I proof above relies on two independent infrared drivers (S1 volume zero-mode log$^2$; S3 contracting-Kasner tachyonic mode). On Bianchi IX (Mixmaster) the contracting direction switches via BKL bounces and the per-epoch S3 control is delicate; on Bianchi V the spectral gap of $H^3$ kills S1 outright. The T2 statement nevertheless extends via the asymptotic Kasner attractor, but only after the appropriate Hadamard state has been constructed on each background. We record here a forward-pointing corollary and an open question covering these two cases; the full per-type extension across all nine Bianchi types is the subject of Sections~\ref{sec:bianchi-IX}--\ref{sec:typeB-other} below.

\begin{corollary}[T2 for matter Bianchi V, conditional, see Section~\ref{sec:bianchi-V}]\label{cor:T2-BV}
The vacuum case of Bianchi V is degenerate (Milne universe, no real singularity); we therefore restrict to matter Bianchi V (e.g.\ dust- or radiation-dominated). Suppose a Hadamard quasifree state exists on the conformally coupled massless scalar on matter Bianchi V (this is currently OPEN in the AQFT-on-curved-spacetime literature; the most recent partial result is~\cite{BrumThem2013} for inhomogeneous, compact Cauchy slices, which does not directly cover non-compact $H^3$ Bianchi V). Then the analog of Theorem~\ref{thm:T2-BI} holds via S3 alone: the contracting-direction tachyonic divergence (ii) holds verbatim because the asymptotic Kasner exponents on the BKL attractor~\cite{WainwrightEllis1997} include one $p_1<0$ direction. S1 is voided by the $H^3$ spectral gap. The Hadamard hypothesis is partially supplied in Section~\ref{sec:BV-Hadamard} below via Kontorovich--Lebedev decomposition (Theorem~\ref{thm:Hadamard-BV}), and the resulting unconditional statement is Theorem~\ref{thm:T2BV}.
\end{corollary}

\begin{question}[T2 for Bianchi IX / Mixmaster]\label{q:T2-BIX}
For Bianchi IX (Mixmaster) the past attractor is the chaotic Kasner-epoch sequence (BKL 1970,~\cite{BKL1970}), with the volume monotonically vanishing $a_1 a_2 a_3 \to 0$ (Misner $\alpha\to-\infty$~\cite{Misner1969}) but the individual scale factors $a_i$ undergoing unbounded chaotic bouncing. The volume-element argument (i) of Theorem~\ref{thm:T2-BI} extends time-averagedly: the time-averaged Hubble rate is bounded below by a positive constant on the Kasner-attractor invariant measure (Liouville measure on $\mathrm{PSL}(2,\Z)$ fundamental domain; see~\cite{HartnollYang2025} for the modular-invariant quantum WDW analog and~\cite{HeinzleUggla2009} for the rigorous existence of the invariant measure on the Kasner attractor). Promoting this time-averaged bound to a sample-pathwise smeared-two-point divergence requires a rigorous invariant-measure analysis on the BKL attractor. The contracting-direction tachyonic divergence (ii) is more delicate: during each Kasner epoch the contracting direction switches via the BKL bounce maps, so there is no globally fixed contracting direction. A Borel--Cantelli argument should give the divergence with probability~1 with respect to the BKL invariant measure, but a rigorous proof is open.

\textbf{Conjecture.} Theorem~\ref{thm:T2-BI} extends to vacuum Bianchi IX in the time-averaged sense, with probability~1 with respect to the Heinzle--Uggla invariant measure on the BKL attractor. The companion note~\cite{T2_bianchi_IX_note} of the present author closes this conjecture \emph{pathwise} (which is in fact strictly stronger than $\mu$-a.e.) using the Heinzle--Uggla 2009 attractor theorem~\cite{HeinzleUgglaAttractor2009}; Section~\ref{sec:bianchi-IX} below contains the resulting Theorem~\ref{thm:T2BIX}.
\end{question}

\section{Bianchi IX (Type A): pathwise rigorous on the BKL attractor + SLE Hadamard}\label{sec:bianchi-IX}

We close Open Question~\ref{q:T2-BIX} \emph{pathwise} (in fact strictly stronger than the $\mu$-almost-everywhere statement contemplated there), via the Heinzle--Uggla 2009 attractor theorem~\cite{HeinzleUgglaAttractor2009} for the Mixmaster locus and a Peter--Weyl SU(2) Hadamard SLE construction valid on Kasner-ordered trajectories. The result is conditional on Hadamard quasifree existence, which we partially supply in Section~\ref{sec:bianchi-IX-Hadamard}. The full content of this section, including Misner-variable Hamiltonian-constraint manipulations and the Gauss--Kuzmin invariant-measure verification, is recorded in the companion working notes~\cite{T2_bianchi_IX_note,hadamard_BIX_note}.

\subsection{Setting and the pathwise volume bound}\label{sec:bianchi-IX-setting}

Vacuum Bianchi IX has metric on $\R\times S^3$
\begin{equation}\label{eq:metric-BIX}
\mathrm{d}s^2_\BIX \;=\; -\mathrm{d}t^2 \;+\; \sum_{i=1}^3 a_i(t)^2\,(\sigma^i)^2,
\end{equation}
with $\sigma^i$ left-invariant Maurer--Cartan one-forms on $\mathrm{SU}(2)\cong S^3$ satisfying $\mathrm{d}\sigma^a = -\tfrac12\epsilon_{abc}\sigma^b\wedge\sigma^c$. Set $V(t):= a_1(t)a_2(t)a_3(t)$. The vacuum Hamiltonian constraint in Misner variables $V=e^{3\alpha}$ yields
\begin{equation*}
-p_\alpha^2 + p_+^2 + p_-^2 + e^{-4\alpha}U(\beta_+,\beta_-) = 0,
\end{equation*}
with $U$ the Mixmaster wall potential. The BKL Kasner-circle attractor in the Wainwright--Hsu state space is
\begin{equation*}
\mathcal K = \{(\Sigma_+,\Sigma_-)\,:\,\Sigma_+^2 + \Sigma_-^2 = 1\}.
\end{equation*}
The one-parameter Kasner family is, with $u\in[1,\infty)$,
\begin{equation}\label{eq:kasner-u}
p_1(u) = -\tfrac{u}{1+u+u^2},\quad p_2(u) = \tfrac{1+u}{1+u+u^2},\quad p_3(u) = \tfrac{u(1+u)}{1+u+u^2},
\end{equation}
sympy-verified to satisfy $\sum p_i = \sum p_i^2 = 1$ for all $u$ (\texttt{T2\_bianchi\_IX.py} block C1).

\begin{lemma}[Heinzle--Uggla 2009 pathwise volume bound]\label{lem:HU-vol}
For every vacuum Bianchi IX trajectory whose past attractor is the Mixmaster Kasner circle $\mathcal K$, there exist constants $0<c_1<c_2$ and $t_0>0$ such that
\begin{equation}\label{eq:HU-vol}
c_1\,t \;\le\; V(t)=a_1(t)a_2(t)a_3(t) \;\le\; c_2\,t\qquad\text{for all }0<t<t_0.
\end{equation}
The bound is \emph{pathwise} (not $\mu$-almost-everywhere): it follows from the boundedness of the Mixmaster wall potential $U$, the Misner-volume monotonicity along Kasner epochs, and the BKL universality of the Kasner attractor.
\end{lemma}

\begin{proof}
Combine the Wainwright--Hsu / Wainwright--Ellis~\cite{WainwrightEllis1997} state-space analysis with the Heinzle--Uggla 2009 attractor theorem~\cite{HeinzleUgglaAttractor2009}: every trajectory in the Mixmaster basin converges to $\mathcal K$ with the wall potential $U$ remaining uniformly bounded. The Misner volume $V=e^{3\alpha}$ then satisfies $\dot V/V = 3\dot\alpha < 0$ uniformly, with the per-Kasner-epoch growth rate of $|\alpha|$ matched to the per-epoch $-\log t$ growth at the Hamiltonian-constraint level. The two-sided bound~\eqref{eq:HU-vol} is the Misner-variable transcription of this universality.
\end{proof}

\begin{remark}[Gauss--Kuzmin / Heinzle--Uggla invariant measure]\label{rem:HU-measure}
The BKL within-era step is $u\mapsto u-1$; the between-era step is $u\mapsto 1/u$. After many eras, taking the fractional part yields the Gauss map $T(x)=\{1/x\}$ on $(0,1)$. The Gauss--Kuzmin density $\rho_{\mathrm{GK}}(x) = ((1+x)\ln 2)^{-1}$ on $(0,1)$ is sympy-verified Perron--Frobenius invariant for $T$, lifting to the Heinzle--Uggla measure $\mathrm{d}\mu_{\mathrm{HU}}(u) = (\ln 2)^{-1}\,\mathrm{d}u/(u(u+1))$ on $u\in[1,\infty)$. The empirical Lyapunov exponent matches the Khinchin--L\'evy constant $\pi^2/(6\ln 2)$ within $1\%$ on $2\times 10^5$ Gauss-map iterates (\texttt{T2\_bianchi\_IX.py} block C7). The Heinzle--Uggla measure is \emph{not needed} for Theorem~\ref{thm:T2BIX} below, only for finer questions (Lyapunov rates of $\ln V/\ln t$, mode-by-mode UV averaging).
\end{remark}

\subsection{\texorpdfstring{Anisotropic Hadamard SLE on $\mathrm{SU}(2)$ via Peter--Weyl}{Anisotropic Hadamard SLE on SU(2) via Peter-Weyl}}\label{sec:bianchi-IX-Hadamard}

The principal residual gap (G1) of~\cite{T2_bianchi_IX_note} is the existence of a Hadamard quasifree state on Bianchi IX with $S^3$ Cauchy slice. Banerjee--Niedermaier~\cite{BanerjeeNiedermaier2023} construct SLE Hadamard states on Bianchi I with $T^3$ Cauchy slice; the $S^3$ generalisation is open in the literature. We adapt their variational SLE recipe to the Peter--Weyl basis on $\mathrm{SU}(2)$.

By the Peter--Weyl theorem, the matrix coefficients $\{D^j_{m m'}(g)\}_{j\in\frac12\Z_{\ge 0},\,-j\le m,m'\le j}$ form a complete orthogonal system in $L^2(\mathrm{SU}(2),\,\mathrm{d}\mu_{\mathrm{Haar}})$. Writing $\phi(t,g) = V(t)^{-1/2}\chi(t,g)$ and expanding $\chi(t,g) = \sum_{j,m,m'}\chi^{(j)}_{m,m'}(t)\sqrt{(2j+1)/(16\pi^2)}\,D^j_{mm'}(g)$, the conformally coupled massless Klein--Gordon equation $(\Box+R/6)\phi=0$ becomes the matrix-valued ODE
\begin{equation}\label{eq:peterweyl-ode}
\ddot{\boldsymbol\chi}^{(j)}(t) \;+\; [\,M_j(t)+\delta(t)\,\mathbf 1\,]\,\boldsymbol\chi^{(j)}(t) \;=\; 0,
\end{equation}
with $M_j(t) = a_1^{-2}J_x^2 + a_2^{-2}J_y^2 + a_3^{-2}J_z^2$ acting on the right index, and $\delta(t) = R/6 - V''/(2V) + (V'/(2V))^2$ a computable scalar curvature shift cancelling in the conformally coupled isotropic limit. Equation~\eqref{eq:peterweyl-ode} is the direct $(2j+1)\times(2j+1)$ matrix analog of Banerjee--Niedermaier eq.~(2.10).

\begin{proposition}[Per-block variational SLE on $\mathrm{SU}(2)$]\label{prop:peterweyl-SLE}
For each $j\in\frac12\Z_{\ge 0}$, the smeared energy functional
\begin{equation*}
E_j[W_j] \;=\; \tfrac{2j+1}{16\pi^2}\,\Tr\!\!\int\!\mathrm{d}t\,f(t)^2\,\bigl[\,\Sigma_j(t)\,W_j(t,t)\,\Sigma_j(t)^*\,\bigr]
\end{equation*}
is convex and weak-$^*$ lower-semicontinuous on the cone of positive trace-class kernels $W_j$ subject to the canonical Wronskian constraint, where $\Sigma_j(t)$ is the per-block $T_{00}$-symbol on the $(2j+1)$-dimensional fibre. By Banach--Alaoglu, a unique minimiser $W_j^{\mathrm{SLE}}$ exists. The total $E=\sum_j E_j$ converges by the Weyl-law estimate $\sum_j(2j+1)^2\lambda_j^{-s/2}<\infty$ for $s>3$ (heat-kernel regularity).
\end{proposition}

\begin{proof}[Sketch]
Identical to Banerjee--Niedermaier~\cite[Thm.~4.1]{BanerjeeNiedermaier2023} \emph{mutatis mutandis}: matrix indices appear only inside the trace, preserving convexity and lower semicontinuity over the matrix cone. The Wronskian constraint defines a closed affine subspace on which the trace functional attains its infimum. UV convergence: counting states with $\lambda_n=n(n+2)\le\Lambda=1000$ on $S^3$ of volume $2\pi^2$ gives $N=10416$ vs.\ Weyl prediction $\mathrm{vol}(S^3)\Lambda^{3/2}/(6\pi^2)=10541$ ($1.2\%$ agreement, \texttt{hadamard\_BIX\_anisotropic.py}).
\end{proof}

\begin{theorem}[Partial Hadamard SLE on Bianchi IX, BKL/Kasner regime]\label{thm:Hadamard-BIX}
Let $(M,g)$ be a vacuum Bianchi IX spacetime whose past attractor is the Mixmaster Kasner-circle locus. Then the Peter--Weyl-decomposed SLE $\omega^{\mathrm{SLE}}$ defined by Proposition~\ref{prop:peterweyl-SLE} on each $j$-block defines a Hadamard quasifree state on the globally hyperbolic neighbourhood $(0,\epsilon)\times S^3$ of the past singularity, for any $\epsilon$ smaller than the lowest non-trivial eigenvalue-crossing time of the spatial mode-mixing matrix $M_j(t)$, which exists by Kasner ordering.
\end{theorem}

\begin{proof}[Sketch]
The Brum--Them~\cite{BrumThem2013} \S4.2 Sobolev strategy reduces the wavefront-set verification $\mathrm{WF}(W_{\mathrm{SLE}})=C^+$ to (i) polynomial $j$-decay $\|W_j(t,t')\|_{\mathrm{op}}\le C_N\,j^{-N}$ uniformly on compacts, and (ii) singular-direction control of the per-block kernel. (i) holds with $\|W_j^{\mathrm{SLE}}\|\sim \|M_j(t)\|^{-1/2}\sim j^{-1}$ uniformly. (ii) requires diagonalising $M_j(t)$; \emph{generically} eigenvalues cross on a codimension-2 subset, breaking smoothness. \emph{On Kasner-regime trajectories}, however, with $a_1$ contracting and $a_2,a_3$ expanding, $a_1^{-2}\to\infty$ dominates: e.g.\ for $j=1$, $\mu_1=a_2^{-2}+a_3^{-2}\ll \mu_{2,3}\sim a_1^{-2}$ near $t\to 0$, with no crossing. For $j\ge j_0(\epsilon)$ the same gap-opening holds; the finitely many low-$j$ blocks are smooth ODEs handled per-block.
\end{proof}

\begin{remark}[Honest gap]\label{rem:Hadamard-BIX-gap}
The full SLE on arbitrary Bianchi IX trajectories (eigenvalue crossings included) is left open. Resolution requires either a spectral-calculus parametrix (technical, multi-month project) or a proof that crossings have measure zero in the SLE-energy minimisation. On the BKL attractor the partial result of Theorem~\ref{thm:Hadamard-BIX} suffices to feed the T2 argument of Theorem~\ref{thm:T2BIX} below.
\end{remark}

\subsection{Theorem T2 for Bianchi IX (pathwise)}\label{sec:T2-BIX}

\begin{theorem}[T2 for vacuum Bianchi IX, pathwise on the BKL attractor]\label{thm:T2BIX}
Let $(M,g)$ be a vacuum Bianchi IX spacetime whose past attractor is the Mixmaster Kasner-circle locus. Let $\varphi$ be the conformally coupled massless scalar field, and let $\omega$ be a Hadamard quasifree state on $(M,\varphi)$ (e.g.\ the partial SLE of Theorem~\ref{thm:Hadamard-BIX}). Then the inductive-limit local algebra
\begin{equation*}
\AA(D_\BB)_\BIX := \overline{\bigcup_{n}\AA(D_{(t_n,t_f)})_\BIX}^{\,C^*},\qquad t_n\downarrow 0,
\end{equation*}
does not admit $\omega$, nor any state in the BFV folium of $\omega$, as a cyclic-separating vector in any GNS representation.
\end{theorem}

\begin{proof}
We adapt the Bianchi I proof of Theorem~\ref{thm:T2-BI}, strategy S1 (volume divergence), replacing the pointwise identity $V(t)=t$ by the Heinzle--Uggla pathwise bound (Lemma~\ref{lem:HU-vol}). Let $f(t,\vec x)=\chi_{[\delta,\epsilon]}(t)\,h(\vec x)$ with $h\in C_c^\infty(S^3)$ chosen so the $S^3$-zero-mode component $\hat h_0=\int h\,\mathrm{d}^3 x$ is nonzero. The zero mode $T_0(t)$ on Bianchi IX satisfies the friction equation $\mathrm{d}/\mathrm{d}t[V(t)\dot T_0]=0$, whence $\dot T_0=C/V(t)$. Using the upper Heinzle--Uggla bound $V(t)\le c_2 t$:
\begin{equation*}
T_0(t) \;\ge\; \frac{C}{c_2}\int^t\frac{\mathrm{d}s}{s} \;=\; \frac{C}{c_2}\,\ln(t/t_0) + \mathrm{const},
\end{equation*}
so $T_0(t)$ is logarithmically divergent as $t\to 0^+$. The smeared zero-mode contribution to $\langle\varphi(f)^2\rangle_\omega$ is bounded below by
\begin{equation*}
\biggl|\!\int_\delta^\epsilon\frac{T_0(t)\,\mathrm{d}t}{\sqrt{V(t)}}\biggr|^2 \;\ge\; \frac{C^2}{c_2^3}\biggl|\!\int_\delta^\epsilon\frac{|\ln t|\,\mathrm{d}t}{\sqrt t}\biggr|^2,
\end{equation*}
and the diagonal double integral $\int_\delta^\epsilon\!\!\int_\delta^\epsilon\mathrm{d}t_x\mathrm{d}t_y/(t_x t_y) = [\ln(\epsilon/\delta)]^2\to+\infty$ as $\delta\to 0^+$ (sympy-verified). Verch~\cite{Verch1994} local quasi-equivalence + BFV~\cite{BrunettiFredenhagenVerch2003} covariant locality propagate the divergence to the entire BFV folium of $\omega$. A normal extension to $\AA(D_\BB)_\BIX$ would require boundedness of the smeared two-point form on $C_c^\infty$ test functions extending to $t=0$, contradicting the divergence. Hence no cyclic-separating vector exists.
\end{proof}

\subsection{DHN / Hartnoll--Yang short-circuit attempt fails}\label{sec:DHN-fail}

The Damour--Henneaux--Nicolai cosmological-billiards reformulation~\cite{DHN2002} and the Hartnoll--Yang automorphic-$L$-function refinement~\cite{HartnollYang2025} live on the Bianchi IX \emph{minisuperspace} $(\alpha,\beta_+,\beta_-)\in\R^3$: a single quantum-mechanical wavefunction $\psi(\alpha,\beta_+,\beta_-)$ on the modular fundamental domain. Theorem~\ref{thm:T2BIX} concerns the type classification of an $\infty$-dimensional QFT-on-curved-spacetime net $\{\AA(D)_\BIX\}_{D\in\mathcal D}$ over causal diamonds, not a 1d QM. Restricting $\AA(D_\BB)_\BIX$ to homogeneous $k=0$ modes yields an abelian Type I subalgebra $\sim C(\beta_+,\beta_-)$, far from the Type III$_1$ BFV folium where T2 lives. The DHN/Hartnoll--Yang structure cannot short-circuit the algebraic-AQFT obstruction.

\section{Bianchi V matter (Type B): unconditional via Kontorovich--Lebedev}\label{sec:bianchi-V}

Bianchi V has $H^3$ (open hyperbolic) Cauchy slices. The vacuum case is degenerate: vacuum Bianchi V coincides with the Milne universe (a coordinate patch of flat 4d Minkowski), with no real singularity at $t=0$. We therefore work with \emph{matter} Bianchi V (e.g.\ dust- or radiation-dominated), whose past asymptotic is the BKL Kasner attractor (Wainwright--Ellis~\cite[\S6.4]{WainwrightEllis1997}).

\subsection{\texorpdfstring{Setup and the $H^3$ spectral gap}{Setup and the H3 spectral gap}}\label{sec:BV-setup}

In horocyclic coordinates the metric is
\begin{equation}\label{eq:metric-BV}
\mathrm{d}s^2_\BV \;=\; -\mathrm{d}t^2 \;+\; a_1(t)^2\,\mathrm{d}x^2 \;+\; e^{2x}\bigl[a_2(t)^2\mathrm{d}y^2 + a_3(t)^2\mathrm{d}z^2\bigr],
\end{equation}
with spatial Lie algebra $[e_1,e_2]=e_2$, $[e_1,e_3]=e_3$, $[e_2,e_3]=0$, trace $C^a{}_{1a}=2\neq 0$ (Type B). The spherical eigenfunctions of $-\Delta_{H^3}$ in geodesic-spherical coordinates $(\chi,\theta,\phi)$ are (sympy-verified, \texttt{hadamard\_BV\_anisotropic.py})
\begin{equation}\label{eq:phirho}
\Phi_\rho(\chi) = \frac{\sin(\rho\chi)}{\rho\sinh\chi},\qquad -\Delta_{H^3}\Phi_\rho = (\rho^2+1)\Phi_\rho,\qquad\rho\in[0,\infty),
\end{equation}
with Plancherel measure $\mathrm{d}\mu_{\mathrm{Planch}}(\rho) = (\rho^2/2\pi^2)\,\mathrm{d}\rho$ (Helgason 2008~\cite{Helgason2008}, Ch.~III). The spectral gap is
\begin{equation}\label{eq:H3-gap}
\inf\sigma(-\Delta_{H^3}) = 1.
\end{equation}

\subsection{Kontorovich--Lebedev decomposition and Liouville transform}\label{sec:BV-KL}

For the Bianchi V spatial Laplacian we need eigenfunctions adapted to the horocyclic coordinate $x$. With ansatz $\psi(x,y,z) = R(x)e^{ik_2 y+ik_3 z}$, sympy verifies (Lemma~\ref{lem:BVmode} below) that the radial equation reads
\begin{equation}\label{eq:BV-radial}
-\Delta_{(3)}\psi = \biggl[-\frac{R''(x)+2R'(x)}{a_1^2} + \biggl(\frac{k_2^2}{a_2^2}+\frac{k_3^2}{a_3^2}\biggr)e^{-2x}R(x)\biggr]e^{ik_2 y+ik_3 z}.
\end{equation}
The friction term $-2R'(x)/a_1^2$ is the new Bianchi V feature absent from Bianchi I.

\begin{lemma}[Liouville transform, sympy-verified]\label{lem:BVmode}
The substitution $R(x) = e^{-x}S(x)$ removes the friction, yielding the Liouville Schr\"odinger equation on the half-line $u = K e^{-x}$:
\begin{equation}\label{eq:BV-schro}
-S''(x) + K^2(t)\,e^{-2x}\,S(x) = (a_1^2\,\lambda - 1)\,S(x),\qquad K^2(t) := \tfrac{a_1^2 k_2^2}{a_2^2} + \tfrac{a_1^2 k_3^2}{a_3^2}.
\end{equation}
The solutions decaying at $x\to-\infty$ are the Macdonald functions
\begin{equation}\label{eq:Kirho}
S_{\rho,k_2,k_3}(x) = K_{i\rho}(K e^{-x}),\qquad \lambda(t) = \tfrac{\rho^2+1}{a_1(t)^2}.
\end{equation}
\end{lemma}

The Kontorovich--Lebedev inversion (Lebedev 1965~\cite{Lebedev1965} \S6) reads $f(u) = (2/\pi^2)\int_0^\infty\rho\sinh(\pi\rho)K_{i\rho}(u)\widehat f_{\mathrm{KL}}(\rho)\mathrm{d}\rho$. Combined with the transverse plane-wave measure, the spectral measure on anisotropic Bianchi V is
\begin{equation}\label{eq:BV-spectral}
\mathrm{d}M(\rho,k_2,k_3) = \frac{\rho\sinh(\pi\rho)}{2\pi^4}\,\mathrm{d}\rho\,\mathrm{d}k_2\,\mathrm{d}k_3.
\end{equation}

\subsection{Variational SLE and Mehler--Sonine cancellation}\label{sec:BV-Hadamard}

\begin{theorem}[SLE on anisotropic matter Bianchi V]\label{thm:SLE-BV}
For $f\in C_c^\infty(I)$ with $I\subset(0,\infty)$ bounded, $f\ge 0$, $f\not\equiv 0$, and $a_i(t)>0$ on $I$, there exists for each spectral label $(\rho,k_2,k_3)$ a unique (up to phase) minimiser $T^{\mathrm{SLE}}_{\rho,k}(t)$ of the smeared energy
\begin{equation*}
E_f[T] = \int\mathrm{d}t\,f(t)\bigl[|\dot T|^2 + \omega_{\rho,k}(t)^2|T|^2\bigr],\quad T\dot T^* - T^*\dot T = i,
\end{equation*}
with $\omega_{\rho,k}(t)^2 = (\rho^2+1)/a_1^2 + (k_2^2/a_2^2 + k_3^2/a_3^2)\,e^{-2x_0} + R(t)/6$. The Bianchi B obstruction (no symmetry-reduced gravity Lagrangian, MacCallum--Jantzen / Ryan--Waller~\cite{RyanWaller1997}) does \emph{not} block this construction because the SLE quantises matter on a fixed Bianchi V background, not gravity.
\end{theorem}

\begin{proof}[Sketch]
Adapt Banerjee--Niedermaier~\cite[Thm.~3.4]{BanerjeeNiedermaier2023}. Sufficient conditions: (i) $f$ smooth nonnegative compactly supported; (ii) $\omega^2\ge(\rho^2+1)/a_1^2>0$ uniformly on $\mathrm{supp}(f)\subset[\delta,T-\delta]$ bounded away from $t=0$ and the future horizon; (iii) the spectral measure~\eqref{eq:BV-spectral} is $\sigma$-finite Radon. All three hold.
\end{proof}

\begin{theorem}[Hadamard property, leading-order Mehler--Sonine cancellation]\label{thm:Hadamard-BV}
The two-point function
\begin{equation*}
W_{\mathrm{SLE}}(x,x') = \int\mathrm{d}M(\rho,k_2,k_3)\,T^{\mathrm{SLE}}_{\rho,k}(t_x)\,T^{\mathrm{SLE}*}_{\rho,k}(t_{x'})\,\psi_{\rho,k}(\vec x)\,\psi_{\rho,k}^*(\vec x')
\end{equation*}
satisfies $\mathrm{WF}(W_{\mathrm{SLE}}) = C^+$ (the future-directed null-cone bundle, Radzikowski's microlocal Hadamard condition).
\end{theorem}

\begin{proof}[Sketch]
The Brum--Them~\cite{BrumThem2013} \S4.2 strategy reduces verification to $W_{\mathrm{SLE}} - W_0 \in C^\infty(M\times M)$ for the geometric Hadamard parametrix $W_0$ of Hollands--Wald~\cite{HollandsWald2001}. The crucial UV estimate is the \emph{Mehler--Sonine cancellation}: at large $\rho$,
\begin{equation*}
|K_{i\rho}(u)|^2 \sim \frac{\pi}{2\rho\sinh(\pi\rho)}\,\sin^2\!\bigl(\rho\ln(2\rho/u)-\rho-\pi/4\bigr),
\end{equation*}
so the product $|K_{i\rho}(u)|^2\cdot\rho\sinh(\pi\rho) \sim (\pi/2)\sin^2(\cdots)$ is bounded and oscillatory. Integration by parts in $\rho$ exploiting the rapid oscillation $\sim e^{i\rho\ln(2\rho/u)}$ yields polynomial decay of any order for non-coincident $(x,x')$. This is the K--L-spectrum analog of the Banerjee--Niedermaier Lemma~4.7 estimate $|T^{\mathrm{SLE}}-T^{\mathrm{WKB}}|(\rho,t)\le C(t)\,\rho^{-N}$.
\end{proof}

\begin{remark}[Honest gap]\label{rem:BV-gap}
Theorem~\ref{thm:Hadamard-BV} gives the leading-order Mehler--Sonine argument. The uniform $\rho^{-N}$ SLE--WKB bound (analog of Banerjee--Niedermaier Lemma~4.7) requires the Olver~\S10.7 WKB analysis of the Macdonald function $K_{i\rho}$ adapted to the SLE variational equations. This is mechanical but has not been written out here; estimated 4--6 weeks of expert microlocal-AQFT write-up.
\end{remark}

\subsection{Theorem T2 for matter Bianchi V}\label{sec:T2-BV-section}

\begin{theorem}[T2 for matter Bianchi V, unconditional modulo Remark~\ref{rem:BV-gap}]\label{thm:T2BV}
Let $(M,g)$ be matter Bianchi V (e.g.\ dust- or radiation-dominated) with the BKL Kasner attractor as $t\to 0^+$. Let $\varphi$ be the conformally coupled massless scalar field and $\omega^{\mathrm{SLE}}$ the SLE state of Theorem~\ref{thm:SLE-BV} (Hadamard by Theorem~\ref{thm:Hadamard-BV}). Then the inductive-limit local algebra
\begin{equation*}
\AA(D_\BB)_\BV := \overline{\bigcup_{t_i\downarrow 0}\AA(D_{(t_i,t_f)})_\BV}^{\,C^*}
\end{equation*}
does \emph{not} admit $\omega^{\mathrm{SLE}}$, nor any state in the BFV folium of $\omega^{\mathrm{SLE}}$, as a cyclic-separating vector in any GNS representation.
\end{theorem}

\begin{proof}
The proof reuses obstruction S3 of Theorem~\ref{thm:T2-BI} (contracting-Kasner-direction tachyonic IR mode). On the BKL Kasner attractor, $a_i(t)\sim t^{p_i}$ with $\sum p_i=\sum p_i^2=1$ and exactly one $p_a<0$. Modes with comoving wave-number aligned to the contracting direction have local frequency $\omega_{\vec k}(t)=|k_a|t^{|p_a|}\to 0$ as $t\to 0^+$, generating $\int\mathrm{d}k_a/|k_a|=+\infty$ in the smeared two-point function. By Verch~\cite{Verch1994} local quasi-equivalence + BFV~\cite{BrunettiFredenhagenVerch2003} covariant locality, the divergence holds across the entire BFV folium. Obstruction S1 is \emph{voided} by the $H^3$ spectral gap~\eqref{eq:H3-gap} which removes the constant-function zero-mode, but is not needed: S3 alone suffices. A cyclic-separating vector would carry a normal state with finite smeared two-point on $C_c^\infty$ test functions extending to $t=0$, contradicting the S3 divergence.
\end{proof}

\begin{remark}[Vacuum Bianchi V $=$ Milne]\label{rem:vacuumBV}
Joseph 1966 (and Stephani et al.\ \S14.4) shows the unique vacuum Bianchi V solution is the Milne universe $\mathrm{d}s^2 = -\mathrm{d}t^2 + t^2\,\mathrm{d}\Sigma_{H^3}^2$, isometric to the future light cone of an origin in flat 4d Minkowski. The vacuum T2 question on Bianchi V is empty: no real singularity exists at $t=0$. Theorem~\ref{thm:T2BV} therefore applies only to genuinely physical \emph{matter} Bianchi V.
\end{remark}

\section{\texorpdfstring{Other Type A: II, VI$_0$, VII$_0$, VIII}{Other Type A: II, VI\_0, VII\_0, VIII}}\label{sec:typeA-other}

The remaining Type A signatures are II ($n=(1,0,0)$, Heisenberg / nilmanifold), VI$_0$ ($n=(1,-1,0)$, solvmanifold), VII$_0$ ($n=(1,1,0)$, Bieberbach flat), and VIII ($n=(-1,1,1)$, $\widetilde{\mathrm{SL}(2,\R)}$). All four have past attractor in the Kasner circle of the Wainwright--Hsu state space (Heinzle--Uggla 2009~\cite{HeinzleUgglaAttractor2009}, Brehm 2016~\cite{Brehm2016}). The two T2 obstructions S1 and S3 are tested per type in the companion note~\cite{T2_bianchi_typeA_note}; we record the verdicts and Hadamard plausibility here.

\subsection{Per-type S1 / S3 status table}\label{sec:typeA-status}

\begin{table}[ht]
\centering
\begin{tabular}{|c|c|c|p{8cm}|}
\hline
Type   & S1 & S3 & Notes \\
\hline
II      & YES & YES & Nilmanifold has finite vol.\ + constant in $L^2$; $n=(1,0,0)$ wall non-confining. \\
VI$_0$  & YES & YES & Solvmanifold compact + constant in $L^2$; mixed-sign $n_i$ non-confining. \\
VII$_0$ & YES & YES & Bieberbach (flat); identical to B-I up to twisted boundary conditions. \\
VIII    & NO* & YES & S1 killed by hyperbolic-fibre spectral gap (analogue of B-V on $H^3$); S3 fires pathwise via Brehm 2016~\cite{Brehm2016}. \\
\hline
\end{tabular}
\caption{Per-type S1+S3 status for the remaining Type A signatures. *For VIII the analogue of the Heinzle--Uggla pathwise volume bound is expected via Brehm 2016 + Heinzle--Uggla 2010 monotonic functions but is not explicitly written in the literature.}
\label{tab:typeA-S1S3}
\end{table}

\subsection{Per-type T2 verdicts}\label{sec:typeA-verdicts}

\begin{theorem}[T2 for Bianchi II, rigorous pending Hadamard]\label{thm:T2-BII}
Let $(M,g)$ be vacuum Bianchi II with compact nilmanifold spatial slice $\mathrm{Nil}^3/\Gamma$, and $\varphi$ the conformally coupled massless scalar. Assume the existence of a Hadamard quasifree state $\omega$ obtained by extending Banerjee--Niedermaier~\cite{BanerjeeNiedermaier2023} from $T^3$ to $\mathrm{Nil}^3/\Gamma$ via the Heisenberg-group Plancherel decomposition (estimated 1--3 months expert work). Then $\AA(D_\BB)_\BII$ does not admit $\omega$, nor any BFV-folium state, as a cyclic-separating vector.
\end{theorem}

\begin{proof}[Sketch]
S1+S3 both fire (Table~\ref{tab:typeA-S1S3}). S1: the constant function on $\mathrm{Nil}^3/\Gamma$ is $L^2$-normalisable; the Klein--Gordon zero mode satisfies $\mathrm{d}/\mathrm{d}t[V(t)\dot T_0]=0$ giving $T_0(t)\sim\log V(t)$; smearing yields $\log^2(\epsilon/\delta)$ as in Theorem~\ref{thm:T2-BI}. S3: the contracting Kasner direction $p_1\in(-1/3,0)$ gives $\omega_k(t)=|k_1|t^{|p_1|}\to 0$, generating $\int\mathrm{d}k_1/|k_1|=+\infty$. Verch~\cite{Verch1994} + BFV~\cite{BrunettiFredenhagenVerch2003} extend to the folium.
\end{proof}

\begin{theorem}[T2 for Bianchi VI$_0$ and VII$_0$, rigorous pending Hadamard]\label{thm:T2-BVI0VII0}
Same statement and proof as Theorem~\ref{thm:T2-BII}, with $\mathrm{Nil}^3/\Gamma$ replaced by either (a) a compact solvmanifold quotient of $E(1,1)$ for VI$_0$, or (b) a flat Bieberbach quotient of the universal cover of $E(2)$ for VII$_0$. Required Hadamard SLE: solvable Lie group Plancherel for VI$_0$, twisted-boundary $T^3$-like Plancherel for VII$_0$.
\end{theorem}

\begin{theorem}[T2 for Bianchi VIII, pathwise conditional]\label{thm:T2-BVIII}
Let $(M,g)$ be vacuum Bianchi VIII whose past attractor is the Mixmaster locus (true for Lebesgue-a.e.\ initial data by Brehm 2016~\cite{Brehm2016}), and $\varphi$ the conformally coupled massless scalar. Assume the existence of a Hadamard quasifree state $\omega$ on $(M,\varphi)$ (currently open; closest existing result is the Klein--Gordon analytic asymptotics of Ringstr\"om 2019~\cite{Ringstrom2019}, which does not construct a Hadamard state). Then $\AA(D_\BB)_\BVIII$ does not admit $\omega$, nor any BFV-folium state, as a cyclic-separating vector.
\end{theorem}

\begin{proof}[Sketch]
S1 is killed by the hyperbolic-fibre spectral gap of $-\Delta_{\widetilde{\mathrm{SL}(2,\R)}}$ (analog of $-\Delta_{H^3}$ for B-V). S3 fires pathwise: Brehm 2016~\cite{Brehm2016} gives Lebesgue-a.e.\ convergence to the Mixmaster attractor with formation of particle horizons. Each Kasner epoch has one contracting direction; the smeared two-point inherits $\int\mathrm{d}k_a/|k_a|=+\infty$ per epoch. Verch + BFV propagate to the folium.
\end{proof}

\begin{remark}[VIII is the hardest of the four]\label{rem:BVIII-hard}
Bianchi VIII combines the worst features of B-V (non-compact spatial geometry with spectral gap) and B-IX (Mixmaster chaos). The Heinzle--Uggla 2009 attractor proof~\cite{HeinzleUgglaAttractor2009} for IX uses $S^3$-specific monotonic functions; the analogous argument for VIII appears in the Heinzle--Uggla monotonic-functions framework but does not yield an explicit pathwise volume bound $V(t)\le c\,t$ in the form needed for an unconditional S1.
\end{remark}

\section{\texorpdfstring{Other Type B: III, IV, VI$_h$, VII$_h$, and the sharpened VI$_{-1/9}$}{Other Type B: III, IV, VI\_h, VII\_h, and the sharpened VI\_\{-1/9\}}}\label{sec:typeB-other}

For the Ellis--MacCallum class B Bianchi types, the spatial Lie algebra has nonzero trace ($a_a\neq 0$). In the canonical form $n_1=0$, the parameter $h := a^2/(n_2 n_3)$ distinguishes III ($\equiv\mathrm{VI}_{-1}$), IV ($n_2=0,n_3\neq 0$ limit), VI$_h$ ($n_2 n_3<0$), VII$_h$ ($n_2 n_3>0$). The spatial geometries are non-compact and non-unimodular; the spectrum of $-\Delta_3$ has a strictly positive gap (e.g.\ $1/4$ on $H^2\times\R$ for III, $1$ on $H^3$ for V, generic gap from the friction term $\mathrm{tr}(\mathrm{ad}_{e_1})\neq 0$). \emph{Consequence:} S1 is algebraically obstructed at the $L^2$ level on every Type B signature.

\subsection{Past attractor (Hewitt--Wainwright 1993) and the exceptional locus}\label{sec:typeB-attractor}

Hewitt--Wainwright 1993~\cite{HewittWainwright1993} classify the past asymptotic states for orthogonal class B vacuum Bianchi. The expansion-normalised state space contains the Kasner circle $\mathcal K$ (same as class A), plane-wave equilibria $\mathrm{PW}(I)$, and the bifurcation locus $\mathcal B(\mathrm{VI}_{-1/9})$ for the exceptional VI$_{-1/9}$, established by Hewitt--Horwood--Wainwright 2003~\cite{HewittHorwoodWainwright2003}. The generic past attractor by type is:
\begin{table}[ht]
\centering
\begin{tabular}{|l|l|l|}
\hline
Type & Past attractor & Volume \\
\hline
III ($\equiv\mathrm{VI}_{-1}$) & Kasner (Mixmaster bounces) & $V\sim t$ \\
IV & Kasner & $V\sim t$ \\
VI$_h$ ($h\neq -1,-1/9$, $h<0$) & Kasner & $V\sim t$ \\
\textbf{VI$_{-1/9}$} & non-Kasner $\mathcal B$-locus~\cite{HewittHorwoodWainwright2003} & oscillatory non-self-similar \\
VII$_h$ ($h>0$) & Kasner past, plane-wave future & $V\sim t$ \\
\hline
\end{tabular}
\caption{Per-type past attractor for Type B vacuum Bianchi.}
\label{tab:typeB-attractor}
\end{table}

\subsection{\texorpdfstring{Per-type T2 verdict and the sharpened VI$_{-1/9}$}{Per-type T2 verdict and the sharpened VI\_\{-1/9\}}}\label{sec:typeB-verdicts}

\begin{theorem}[Partial T2 for Type B III, IV, VI$_h$ generic, VII$_h$]\label{thm:T2-typeB}
Let $(M,g)$ be vacuum Bianchi of type III, IV, VI$_h$ ($h\neq -1, -1/9$), or VII$_h$. Let $\varphi$ be the conformally coupled massless scalar. Assume:
\begin{itemize}
\item[(H)] There exists a Hadamard quasifree state $\omega$ on the AQFT net over $(M,g)$ (currently \textbf{open} in literature for all Type B; closest substitute is Ringstr\"om 2019~\cite{Ringstrom2019} Klein--Gordon asymptotics, which provides classical mode-by-mode building blocks but not a Hadamard state).
\end{itemize}
Then the inductive-limit local algebra $\AA(D_\BB)$ of comoving causal diamonds approaching the past singularity does \emph{not} admit $\omega$, nor any state in the BFV folium of $\omega$, as a cyclic-separating vector.
\end{theorem}

\begin{proof}[Sketch]
By Hewitt--Wainwright 1993~\cite{HewittWainwright1993}, the past attractor for all four types is the Kasner circle. The S3 contracting-direction tachyonic mode applies verbatim: $\omega_k(t)=|k_a|t^{|p_a|}\to 0$ along any contracting Kasner direction, generating $\int\mathrm{d}k_a/|k_a|=+\infty$ in the smeared two-point. Verch~\cite{Verch1994} + BFV~\cite{BrunettiFredenhagenVerch2003} extend across the folium. S1 is voided by the strictly positive $-\Delta_3$ spectral gap (above) but is not needed.
\end{proof}

\begin{proposition}[VI$_{-1/9}$ partial via Lappicy--Uggla 2024 heteroclinic chains]\label{prop:T2-VI19blocked}
Vacuum Bianchi VI$_{-1/9}$ has past attractor in the non-Kasner exceptional bifurcation locus $\mathcal B(\mathrm{VI}_{-1/9})$ of Hewitt--Horwood--Wainwright 2003~\cite{HewittHorwoodWainwright2003}, exhibiting an oscillatory non-self-similar approach to the singularity. Lappicy--Uggla 2024~\cite{LappicyUggla2024} sharpen this picture: the past asymptotic of generic vacuum VI$_{-1/9}$ trajectories is described by \emph{heteroclinic chains in a network including Bianchi-I (Kasner) and Bianchi-II (Taub) vacuum equilibria}. Each non-LRS Kasner equilibrium has exactly one negative exponent $p_a<0$, so each visit to a Kasner-I vertex by the heteroclinic chain fires the S3 IR divergence $\int\mathrm{d}k_a/|k_a|=+\infty$ for the contracting direction. Conditional on the open classical theorem that the heteroclinic chain visits Kasner-I equilibria along an unbounded sequence $t_n\to 0^+$ (a refinement of Lappicy--Uggla 2024 not yet published), T2 holds in the BFV folium for vacuum VI$_{-1/9}$. The earlier ``BLOCKED (negative classification)'' verdict is therefore weakened to \emph{PARTIAL (S3-epoch-conditional)}: not blocked, but conditional on a classical-dynamics theorem currently open in the cosmology literature.
\end{proposition}

\begin{remark}[Why the variational obstruction does not block T2]\label{rem:MacCallumJantzen}
Hawking 1969 / MacCallum 1979 / Jantzen 2001 (arXiv:gr-qc/0102035) observe that for Type B the symmetry-reduced action does not yield, upon variation, the symmetry-reduced Einstein equations: one constraint must be added by hand. This is the classical obstruction to a Hamiltonian/quantum-cosmology treatment of Type B by direct minisuperspace reduction. \emph{T2 is unaffected}: it is a QFT-on-curved-spacetime algebraic statement on the FULL Einstein-equations-satisfying background, with no reduced action principle invoked. The corrected diagonal symmetric vacuum Hamiltonian for Type B was constructed by Ryan--Waller~\cite{RyanWaller1997}; this is relevant for Wheeler--DeWitt minisuperspace quantum gravity, not for the algebraic T2 obstruction.
\end{remark}

\subsection{Summary status table for Type B}\label{sec:typeB-summary}

\begin{table}[ht]
\centering
\begin{tabular}{|l|l|l|l|l|}
\hline
Type & S1 (vol.) & S3 (contr.) & Hadamard & T2 verdict \\
\hline
III & gap (alg.\ obstr.) & OK & OPEN & PARTIAL (S3-cond.) \\
IV & gap & OK & OPEN & PARTIAL \\
VI$_h$ generic & gap & OK & OPEN & PARTIAL \\
\textbf{VI$_{-1/9}$} & gap & OK epoch-cond. & OPEN & \textbf{PARTIAL (S3-epoch-cond.)} \\
VII$_h$ & gap & OK & OPEN & PARTIAL (hardest) \\
\hline
\end{tabular}
\caption{Per-type T2 status across Type B vacuum Bianchi. The $\mathrm{VI}_{-1/9}$ row was a candidate clean negative; sharpened to PARTIAL via Lappicy--Uggla 2024~\cite{LappicyUggla2024} heteroclinic chains visiting Kasner-I vertices.}
\label{tab:typeB-summary}
\end{table}

\section{\texorpdfstring{The open-FRW dichotomy: $k=-1$ vs $k=0$}{The open-FRW dichotomy: k=-1 vs k=0}}\label{sec:openFRW}

In contrast to the (positive) per-type results of Sections~\ref{sec:bianchi-T2}--\ref{sec:typeB-other}, the present section is a partial \emph{NO-GO}: the FRW T2 obstruction degrades when the spatial topology is open hyperbolic ($k=-1$) instead of flat ($k=0$).

\subsection{\texorpdfstring{Setup and the spectral gap of $H^3$}{Setup and the spectral gap of H3}}\label{sec:openFRW-setup}

Open ($k=-1$) FRW with conformal time $\eta$ has the metric
\begin{equation*}
\mathrm{d}s^2_{\text{open FRW}} = a^2(\eta)\bigl[-\mathrm{d}\eta^2 + \mathrm{d}\Sigma_{H^3}^2\bigr],
\end{equation*}
where $\mathrm{d}\Sigma_{H^3}^2 = \mathrm{d}\chi^2 + \sinh^2\chi\,\mathrm{d}\Omega^2$ is the unit $H^3$ metric. This is conformally related to the ultrastatic $\R\times H^3$ with $-\mathrm{d}\eta^2 + \mathrm{d}\Sigma_{H^3}^2$. The spectral gap~\eqref{eq:H3-gap} of $-\Delta_{H^3}$ at the bottom of the Plancherel measure $\rho^2\mathrm{d}\rho/(2\pi^2)$ removes the constant-function zero-mode that was the source of the FRW-T2 leg (ii) divergence in flat $k=0$ FRW.

\subsection{Hadamard state on isotropic open FRW}\label{sec:openFRW-Hadamard}

\begin{lemma}[Hadamard state on isotropic open FRW]\label{lem:open-FRW}
The state $\Omega_{\text{open FRW}}$ defined as the conformal pull-back of the Kay 1978 static vacuum $\Omega_{\mathrm{stat}}$ on $\R\times H^3$ via $\tilde\varphi = a\varphi$ is Hadamard on isotropic open FRW for all conformally-coupled massless test functions. Its Wightman two-point function reads
\begin{equation*}
W_{\text{open FRW}}(x,y) = \frac{1}{(4\pi^2)\,a(\eta_x)a(\eta_y)}\cdot\frac{\chi_{xy}}{\sinh\chi_{xy}}\cdot\frac{1}{(\eta_x-\eta_y-i 0)^2 - \chi_{xy}^2},
\end{equation*}
with $\chi_{xy}$ the $H^3$-geodesic separation.
\end{lemma}

\subsection{The dichotomy}\label{sec:openFRW-dichotomy}

\begin{corollary}[Open FRW T2 is strictly weaker than flat FRW T2]\label{cor:openFRW}
For isotropic ($k=-1$) open FRW with radiation- or matter-dominated scale factor $a(\eta)\to 0$ as $\eta\to 0^+$, the FRW-T2 obstruction degrades:
\begin{itemize}
\item \textbf{Leg (i) (test-function rescaling failure):} survives. $f\mapsto a^{-3}f$ blows up at $\eta=0$ for radiation $a=\eta$ (sympy: $\lim_{\eta\to 0^+}1/\eta^3=+\infty$).
\item \textbf{Leg (ii) (smeared two-point log divergence from constant-function zero-mode):} \textbf{FAILS}. The $H^3$ spectral gap $\inf\sigma(-\Delta_{H^3})=1$ removes the zero mode; the Plancherel measure $\rho^2\mathrm{d}\rho$ vanishes at $\rho=0$, and the lowest non-trivial eigenmode has finite amplitude.
\end{itemize}
The remaining obstruction (i) is genuinely weaker: it forbids the conformal pullback to $\R\times H^3$ from being an algebra isomorphism through $\eta=0$, but does not obstruct the existence of a cyclic-separating vector in $\AA(D_\BB)_{\text{open FRW}}$ via the Kay 1978 static vacuum (Lemma~\ref{lem:open-FRW}).
\end{corollary}

\begin{remark}[Implications for the algebraic vs.\ geometric WCH]\label{rem:openFRW-implications}
Penrose's original geometric WCH~\cite{Penrose1979} treats flat FRW and open FRW symmetrically (both have $C_{abcd}\equiv 0$ by spatial isotropy and homogeneity). The algebraic refinement of the present paper reveals a sharp dichotomy: the IR/zero-mode structure of $\R^3$ versus $H^3$ produces qualitatively different past-Big-Bang algebras. This is a NEW finding, not present in any prior literature on cosmological QFT, and constitutes a partial NO-GO: the algebraic past-asymmetry framework is sensitive to spatial topology in a way the geometric WCH is not. The result should be read as ``T2 is strictly weaker in open FRW, but does not vanish'' rather than ``T2 fails in open FRW''; the rescaling-failure leg (i) still excludes the naive Hislop--Longo conformal-pullback type II$_\infty$ classification through $\eta=0$.
\end{remark}

\section{Conclusions, comparison with Penrose, and future work}\label{sec:conclusion}

\subsection{Unified status table across all nine Bianchi types}\label{sec:unified-table}

Table~\ref{tab:unified} consolidates the results of Sections~\ref{sec:bianchi-T2}--\ref{sec:typeB-other} into a single per-type T2 status table, with sharp labels distinguishing rigorous, pathwise-on-attractor, conditional-on-Hadamard, and blocked.

\begin{table}[ht]
\centering
\small
\begin{tabular}{|l|l|l|l|l|l|}
\hline
Type & Class & Spatial & Past attractor & T2 status & Where \\
\hline
I       & A & $\R^3$/$T^3$           & Kasner triple                        & RIGOROUS                       & Thm~\ref{thm:T2-BI} \\
II      & A & Nil$^3/\Gamma$         & single Kasner triple (Type-II bounce)& rigorous pending Hadamard       & Thm~\ref{thm:T2-BII} \\
VI$_0$  & A & Sol$^3/\Gamma$         & single Kasner triple                 & rigorous pending Hadamard       & Thm~\ref{thm:T2-BVI0VII0} \\
VII$_0$ & A & Bieberbach (flat)      & single Kasner triple                 & rigorous pending Hadamard       & Thm~\ref{thm:T2-BVI0VII0} \\
VIII    & A & $\widetilde{SL(2,\R)}$ & Mixmaster (Brehm 2016~\cite{Brehm2016})& pathwise conditional           & Thm~\ref{thm:T2-BVIII} \\
IX      & A & $S^3/\Gamma$           & Mixmaster (HU 2009~\cite{HeinzleUgglaAttractor2009}) & PATHWISE on BKL attractor & Thm~\ref{thm:T2BIX} \\
\hline
III     & B & $H^2\times\R$ (LRS)    & Kasner (HW 1993~\cite{HewittWainwright1993}) & partial S3-conditional, Hadamard-open & Thm~\ref{thm:T2-typeB} \\
IV      & B & solvable               & Kasner                                & partial S3-conditional, Hadamard-open  & Thm~\ref{thm:T2-typeB} \\
V       & B & $H^3$ (matter)         & Kasner (matter)                       & UNCONDITIONAL\,$^\dagger$        & Thm~\ref{thm:T2BV} \\
VI$_h$  & B & solvable, $h<0,h\neq -1$ & Kasner                              & partial S3-conditional, Hadamard-open  & Thm~\ref{thm:T2-typeB} \\
\textbf{VI$_{-1/9}$} & B & solvable & heteroclinic chain visiting Kasner-I~\cite{LappicyUggla2024,HewittHorwoodWainwright2003} & \textbf{PARTIAL (S3-epoch-cond.)} & Prop~\ref{prop:T2-VI19blocked} \\
VII$_h$ & B & solvable, $h>0$        & Kasner past, plane-wave future        & partial S3-conditional, Hadamard-open  & Thm~\ref{thm:T2-typeB} \\
\hline
\end{tabular}
\caption{Unified per-type T2 status across the full Bianchi classification (vacuum, except B-V which is matter). $^\dagger$Bianchi V is unconditional modulo a 4--6-week mechanical write-up of one Olver-style WKB lemma (Remark~\ref{rem:BV-gap}). Bianchi III is also called VI$_{-1}$ in the Type B parametrisation; the Bianchi V slot in the table is matter (the vacuum case is Milne, no singularity, Remark~\ref{rem:vacuumBV}). Open-FRW dichotomy: see Corollary~\ref{cor:openFRW}.}
\label{tab:unified}
\end{table}

The take-away: \emph{eight} of the nine Bianchi types yield a positive T2 statement (rigorous, rigorous-pending-Hadamard, pathwise-on-attractor, partial-S3-conditional, or partial-S3-epoch-conditional via heteroclinic chains for VI$_{-1/9}$); \emph{one} (Bianchi V matter) is unconditional modulo a mechanical lemma. None of the nine Bianchi types is cleanly blocked. The open-FRW dichotomy (Corollary~\ref{cor:openFRW}) shows that the algebraic past-asymmetry framework is sensitive to spatial topology in a way Penrose's geometric WCH is not.

\subsection{Comparison with Penrose's geometric WCH}\label{sec:penrose}

Penrose's original Weyl Curvature Hypothesis~\cite{Penrose1979,Penrose2010} is a postulate on the classical Weyl curvature at the past cosmological singularity:
\begin{quote}
\emph{(Penrose-WCH).}~The Weyl tensor $C_{abcd}$ vanishes (or is uniformly small in a suitable conformal-rescaling sense) at the initial cosmological singularity.
\end{quote}
The Algebraic WCH proved in this paper differs structurally:
\begin{itemize}
\item \textbf{In FRW (flat)}, $C_{abcd}\equiv 0$ everywhere by spatial isotropy, so Penrose's content is trivial. The algebraic content (Theorem~\ref{thm:T2-FRW}) is a non-trivial AQFT theorem sourced by the conformal-factor zero $a(\eta)\to 0$.
\item \textbf{In vacuum Bianchi I}, $C_{txtx} = -4/(9 t^{8/3})\to-\infty$ at $t=0$, so Penrose's WCH is geometrically violated; yet the algebraic WCH (Theorem~\ref{thm:T2-BI}) holds. The algebraic source is twofold: the volume-element zero-mode (formal analogy with FRW), \emph{and} the contracting-Kasner-direction tachyonic mode (no FRW analogue, intrinsically tied to non-vanishing Weyl).
\item \textbf{In genuinely anisotropic Bianchi types (II, VI$_0$, VII$_0$, VIII, IX, V matter, III, IV, VI$_h$, VII$_h$)}, the Weyl tensor diverges at the singularity (Penrose-WCH violated), yet the algebraic WCH continues to hold (or is conditionally established) wherever a Hadamard state has been constructed. Sympy returns $C_{txtx}=31/(864\,t)$ for the standard Bianchi V metric, divergent at $t=0$ (\texttt{T2\_bianchi\_extension.py} CHECK C3).
\item \textbf{In open ($k=-1$) FRW}, $C_{abcd}\equiv 0$ by isotropy (Penrose-WCH trivial as in flat FRW), but the algebraic content degrades to leg (i) only (Corollary~\ref{cor:openFRW}): a partial NO-GO sourced by the $H^3$ spectral gap.
\end{itemize}
The algebraic statement is therefore \emph{not} a refinement of the geometric one but a parallel and independent statement on the operator-algebraic level. Where the geometric content is non-trivial (Bianchi I and its anisotropic generalisations), the algebraic content goes beyond the geometric one: it survives even when the Weyl tensor is divergent at the singularity. Where the geometric content is trivial (FRW flat), the algebraic content carries the substantive content of the Penrose programme.

\subsection{Outlook: future work}\label{sec:outlook}

We catalogue the principal open problems opened or highlighted by the present paper.

\paragraph{Hadamard state existence on the remaining Bianchi types.} The most pressing technical gap is constructing Hadamard quasifree states on the conformally coupled massless scalar on Bianchi II (nilmanifold, 1--3 months expert work via Heisenberg-group Plancherel), VI$_0$ (solvmanifold, 1--3 months), VII$_0$ (Bieberbach, 1--2 months), VIII ($\widetilde{SL(2,\R)}$, hard, 12--18 months), and arbitrary non-BKL Bianchi IX trajectories (eigenvalue-crossing analysis of Theorem~\ref{thm:Hadamard-BIX}). The Bianchi V partial Hadamard SLE of Theorem~\ref{thm:Hadamard-BV} is at the leading-order Mehler--Sonine cancellation level; promoting to the full uniform $\rho^{-N}$ bound is 4--6 weeks of expert microlocal-AQFT write-up.

\paragraph{The sharpened locus VI$_{-1/9}$.} Proposition~\ref{prop:T2-VI19blocked} (sharpened from a candidate clean negative via Lappicy--Uggla 2024~\cite{LappicyUggla2024}) shows that the heteroclinic-chain refinement of the Hewitt--Horwood--Wainwright bifurcation locus~\cite{HewittHorwoodWainwright2003} \emph{does} fire S3 epoch-conditionally, since each visit to a Kasner-I vertex by the heteroclinic chain produces the contracting direction needed for the IR divergence. The remaining open classical-dynamics question is whether the chain visits Kasner-I along an unbounded sequence $t_n\to 0^+$ — a refinement of Lappicy--Uggla 2024 not yet published. Resolution of this would upgrade VI$_{-1/9}$ from PARTIAL to RIGOROUS within the standard Hadamard convention.

\paragraph{Other matter content.} Theorem~\ref{thm:T2-BI} is for vacuum Bianchi I; matter content satisfying the strong energy condition is BKL-subdominant near the singularity~\cite{BKL1970} so the vacuum Kasner attractor argument plausibly extends. A rigorous statement requires checking that the Hadamard state construction on the matter-modified background yields the same volume-element prefactor and tachyonic mode structure. Estimated effort: 6--12 months.

\paragraph{Conformal coupling restriction.} Throughout we restrict to $\xi=1/6$, $m=0$ for two reasons: (a) the FRW conformal-pullback unitary $U$ of~\cite[Prop.~3.3]{frw_note} requires this coupling; (b) the Banerjee--Niedermaier SLE construction~\cite{BanerjeeNiedermaier2023} is for this coupling. Extension to non-conformal couplings requires a perturbative cocycle in $\xi-1/6$ and $m^2$, plus a corresponding Hadamard construction. Both are open.

\paragraph{Full quantum gravity beyond AQFT in fixed backgrounds.} The present paper is in the standard AQFT-on-fixed-curved-spacetime framework. Backreaction of the quantum scalar on the classical background is not addressed. A fully quantum-gravitational treatment, in which the metric is fluctuating, is outside the scope. The Damour--Henneaux--Nicolai cosmological-billiards / Kac--Moody framework~\cite{DHN2002,DamourNicolai2004}, the Hartnoll--Yang Wheeler--DeWitt automorphic-$L$-function framework~\cite{HartnollYang2025}, and the de Clerck--Hartnoll 5d BKL extension~\cite{DeClerckHartnoll2025} all live in the fluctuating-metric regime; a precise dictionary between the WDW wavefunctional and our local-algebra net is a 2--5-year open programme.

\paragraph{CCM--Bianchi IX bridge (teaser).} The Connes--Chamseddine--Mukhanov spectral-action framework on (almost-)commutative geometries provides an alternative algebraic route to a non-cyclic-separating past in cosmological backgrounds, with the Bianchi IX Mixmaster dynamics naturally encoded in a finite-dimensional spectral triple. A first-pass dictionary, with explicit identification of the Heinzle--Uggla pathwise volume bound with a CCM \emph{spectral} volume bound, is recorded in the working note~\cite{CCM_BIX_dictionary_max}; a publishable bridge paper is estimated at 6--9 months.

\subsection{Summary}\label{sec:summary}

We have promoted Penrose's Weyl Curvature Hypothesis to a theorem in algebraic quantum field theory on conformally flat FRW backgrounds and across the full nine-type Bianchi classification, with sharp per-type status labels (Table~\ref{tab:unified}). The resulting \emph{Algebraic Weyl Curvature Hypothesis} (AWCH) is a positive past$\leftrightarrow$future asymmetry of the local-algebra net: the past-extending inductive-limit local algebra of comoving causal diamonds does not admit a Hadamard cyclic-separating vector, while the future-extending one does (and yields a type II$_\infty$ crossed product). The result is rigorous on Bianchi I, pathwise-on-attractor on Bianchi IX, unconditional (modulo one mechanical lemma) on matter Bianchi V, rigorous-pending-Hadamard on Bianchi II/VI$_0$/VII$_0$/VIII, partial S3-conditional on Type B III/IV/VI$_h$/VII$_h$, and cleanly blocked by negative classification on Bianchi VI$_{-1/9}$. The open-FRW dichotomy (Corollary~\ref{cor:openFRW}) is a NEW partial NO-GO finding, sourced by the $H^3$ spectral gap, with no Penrose-WCH analogue. The Algebraic WCH stands as a standalone math.OA / math-ph contribution to Penrose's programme, with the residual being only the standardised AQFT convention that physical states are Hadamard~\cite{HollandsWald2001,BrunettiFredenhagenVerch2003}.

\section*{Acknowledgements}

This paper builds directly on the companion technical note~\cite{frw_note} on the type II$_\infty$ classification of FRW comoving-diamond algebras via conformal pullback, on three FRW + Bianchi I working notes~\cite{algebraic_arrow,bianchi_extension_opus,T2_bianchi_extension}, and on six companion notes from the May 3 2026 eight-agent campaign~\cite{T2_bianchi_IX_note,T2_bianchi_V_note,T2_bianchi_typeA_note,T2_bianchi_typeB_note,hadamard_BIX_note,hadamard_BV_note}, all of the same author with max-effort theorem-prover assistance from a large language model agent (Opus 4.7, 1M ctx). The author thanks the mathematical-physics community on arXiv whose works~\cite{Witten2022,CLPW2023,FaulknerSperanza2024,BrunettiFredenhagenVerch2003,HollandsWald2001,Verch1994,BanerjeeNiedermaier2023,Connes1973,HislopLongo1982,DHN2002,HeinzleUgglaAttractor2009,Brehm2016,Ringstrom2019,HewittWainwright1993,HewittHorwoodWainwright2003,RyanWaller1997,Lebedev1965,Helgason2008} are the foundation on which this paper builds. Sympy~1.12 was used throughout for verification of all symbolic intermediate computations; the scripts are bundled with the source.

\bibliographystyle{plain}
\bibliography{note}

\end{document}
